% SGA 4 第二册；PDF 1—33 页，止于第 V 讲第 5 节标题前。
\SourcePage{1}
\begin{center}
M. Artin、A. Grothendieck、J.-L. Verdier\\
P. Deligne、B. Saint-Donat 参与\\[1em]
{\Large \Term{topos}{拓扑斯}理论与\Term{scheme}{概形}的\Term{etale-cohomology}{平展上同调}}\\
第二册
\end{center}
\SourcePage{2}
M. Artin、A. Grothendieck、J.-L. Verdier，P. Deligne、B. Saint-Donat 参与。
\SourcePage{3}
\begin{center}
\textbf{\Term{topos}{拓扑斯}理论与\Term{scheme}{概形}的\Term{etale-cohomology}{平展上同调}——第二册}\\
M. Artin、A. Grothendieck、J.-L. Verdier\\
P. Deligne、B. Saint-Donat 参与
\end{center}
\textbf{摘要。}——

\textbf{英文摘要。}——待补。
\SourcePage{4}
% 原页为空白。
\SourcePage{5}
\textbf{原书目录（中文重建；下列数字为原书印刷页码）。}

第 V 讲：\Term{topos}{拓扑斯}中的\Term{cohomology}{上同调}，J.-L. Verdier，1。
引言，1；0. \Term{abelian-category}{阿贝尔范畴}的一般性质，2；1. \Term{flat-module}{平坦模}，5；2. \Term{cech-cohomology}{Čech 上同调}；\Term{cohomology}{上同调}记号，13；3. 与\Term{covering}{覆盖}有关的 Cartan--\Term{leray-spectral-sequence}{Leray 谱序列}，19；4. \Term{acyclic-sheaf}{无环层}，22；5. $R^qu_*$ 及与\Term{topos-morphism}{拓扑斯态射}有关的 Cartan--\Term{leray-spectral-sequence}{Leray 谱序列}，27；6. \Term{local-scheme}{局部} Ext 与\Term{cohomology-with-support}{带支集上同调}，29；7. 附录：\Term{cech-cohomology}{Čech 上同调}，37；8. 附录：\Term{local-inductive-limit}{局部归纳极限}（P. Deligne），49；参考文献，63。

第 V bis 讲：\Term{cohomological-descent}{上同调下降}技术，B. Saint-Donat，63。
引言，63；1. 预备知识，65；2. \Term{cohomological-descent}{上同调下降}方法，74；3. \Term{descent}{下降}判据，88；4. 例子，102；5. 应用，107；参考文献，116。

第 VI 讲：\Term{finiteness-condition}{有限性条件}；\Term{fibered-topos}{纤维拓扑斯}与\Term{fibered-site}{纤维景}；在取\Term{limit}{极限}技术中的应用，A. Grothendieck、J.-L. Verdier，119。
0. 引言，119；1. \Term{topos}{拓扑斯}中\Term{object}{对象}与\Term{arrow}{箭头}的\Term{finiteness-condition}{有限性条件}，121；2. \Term{topos}{拓扑斯}的\Term{finiteness-condition}{有限性条件}，137；3. \Term{topos-morphism}{拓扑斯态射}的\Term{finiteness-condition}{有限性条件}，150；4. 由\Term{gluing}{黏合}得到的\Term{topos}{拓扑斯}中的\Term{finiteness-condition}{有限性条件}，155；5. \Term{functor}{函子} $H^i(X,-)$ 与\Term{filtered-colimit}{滤过归纳极限}的交换，167；6. \Term{fibre-category}{纤维范畴}的\Term{inductive-limit}{归纳极限}与\Term{projective-limit}{逆极限}，169。
\SourcePage{6}
7. \Term{fibered-topos}{纤维拓扑斯}与\Term{fibered-site}{纤维景}，178；8. \Term{fibered-topos}{纤维拓扑斯}的\Term{projective-limit}{逆极限}，197；9. 附录：点存在的判据，227；参考文献，229。

第 VII 讲：\Term{scheme}{概形}的\Term{etale-site}{平展景}与\Term{etale-topos}{平展拓扑斯}，A. Grothendieck，231。
1. \Term{etale-topology}{平展拓扑}，231；2. 层的例子，234；3. \Term{etale-topos}{平展拓扑斯}的\Term{generator}{生成元}；层的 $\varinjlim$ 的\Term{cohomology}{上同调}，236；4. 与其他\Term{topology}{拓扑}的比较，237；5. \Term{scheme}{概形}的\Term{projective-limit}{逆极限}的\Term{cohomology}{上同调}，240。

第 VIII 讲：\Term{fibre-functor}{纤维函子}、\Term{support}{支集}、\Term{finite-morphism}{有限态射}的\Term{cohomology}{上同调}研究，A. Grothendieck，247。
1. \Term{etale-topos}{平展拓扑斯}的\Term{topology}{拓扑}不变性，247；2. 域的谱上的层，249；3. 与\Term{scheme}{概形}的\Term{geometric-point}{几何点}有关的\Term{fibre-functor}{纤维函子}，250；4. \Term{strictly-local-ring}{严格局部环}与\Term{strictly-local-scheme}{严格局部概形}，255；5. 在计算 $R^qf_*$ 的\Term{fiber}{纤维}中的应用，258；6. \Term{support}{支集}，261；7. \Term{fibre-functor}{纤维函子}的\Term{specialization-morphism}{特殊化态射}，263；8. \Term{integral-morphism}{整态射}的两个\Term{spectral-sequence}{谱序列}，269；9. \Term{etale-sheaf}{平展层}的\Term{descent}{下降}，272。
术语索引，277；记号索引，281。

\SourcePage{7}
\SGAchapter{V}{拓扑斯中的上同调}{Cohomologie dans les topos}
\begin{center}J.-L. Verdier\end{center}
\textbf{本讲原目录（中文重建；数字为原书印刷页码）。}
引言，1；0. \Term{abelian-category}{阿贝尔范畴}的一般性质，2；1. \Term{flat-module}{平坦模}，5；2. \Term{cech-cohomology}{Čech 上同调}；\Term{cohomology}{上同调}记号，13；3. 与\Term{covering}{覆盖}有关的 Cartan--\Term{leray-spectral-sequence}{Leray 谱序列}，19；4. \Term{acyclic-sheaf}{无环层}，22；5. $R^qu_*$ 及与\Term{topos-morphism}{拓扑斯态射}有关的 Cartan--\Term{leray-spectral-sequence}{Leray 谱序列}，27；6. \Term{local-scheme}{局部} Ext 与\Term{cohomology-with-support}{带支集上同调}，29；7. 附录：\Term{cech-cohomology}{Čech 上同调}，37；8. 附录：\Term{local-inductive-limit}{局部归纳极限}（P. Deligne），49；参考文献，63。

\textbf{引言。}
本讲介绍\Term{topos}{拓扑斯}的交换的、初等的\Term{cohomological-invariant}{上同调不变量}。第 1 节研究\Term{flat-module}{平坦模}和\Term{ringed-topos}{赋环拓扑斯}的\Term{flat-morphism}{平坦态射}。证明使用一个在实际中通常成立的假设，即\Term{topos}{拓扑斯}有\Term{enough-points}{足够多的点}（第 IV 讲 6）。附录第 8 节在一般情形下重新给出这些证明；Deligne 利用\Term{local-inductive-limit}{局部归纳极限}技术，还把 D. Lazard 关于\Term{flat-module}{平坦模}结构的定理推广到\Term{topos}{拓扑斯}情形。本讲的定理是联系各种\Term{cohomological-invariant}{上同调不变量}的\Term{spectral-sequence}{谱序列}的\Term{existence-theorem}{存在定理}（第 3、5、6 节）。众所周知，即使对于\Term{topological-space}{拓扑空间}，\Term{cech-cohomology}{Čech 上同调}通常也不等于\Term{sheaf-cohomology}{层上同调} [11]。
\SourcePage{8}
附录第 7 节引入一种修改后的 Čech 计算，借助\Term{covering-recouvrement}{覆盖}来得到任意\Term{topos}{拓扑斯}中的\Term{sheaf-cohomology}{层上同调}。在该附录中，我们将使用单纯\Term{covering-recouvrement}{覆盖}（\Term{hypercovering}{超覆盖}），其\Term{homotopy}{同伦}\Term{invariants}{不变量}已在 [1] 中研究（另见 [17]）。

这里引入的\Term{cohomological-invariant}{上同调不变量}之所以是初等的，是因为我们不使用\Term{derived-category}{导出范畴} [12]。熟悉这一语言的读者可以立即用它改述本讲的各个命题，进而把它们推广到\Term{complex}{复形}与\Term{hypercohomology}{超上同调}。本讨论班的后续部分将使用\Term{derived-category}{导出范畴}语言。

这里只讨论交换\Term{cohomology}{上同调}。本讨论班中使用的非交换 $H^1$ 以及非交换 $H^2$，请参阅 [9]。我们取\Term{derived-functor}{导出函子}的那些\Term{functor}{函子}都是\Term{additive}{加性}的；这里不讨论乘法结构（参见 [7]）。

\SGAsection{0}{阿贝尔范畴的一般性质}
本节回顾几个引理，其中大部分见于 [11]。
\begin{SGAresult}{命题}{0.1}
设 $\mathcal A$ 是具有\Term{generator}{生成元}的\Term{abelian-category}{阿贝尔范畴}。下列条件\Term{equivalence}{等价}：

\noindent i) $\mathcal A$ 满足公理 AB 5)：小\Term{direct-sum}{直和}\Term{representable}{可表}；若 $(X_i)_{i\in I}$ 是 $\mathcal A$ 中\Term{object}{对象} $X$ 的一族小的\Term{filtered}{滤过}递增\Term{subobject}{子对象}，$Y$ 是 $X$ 的\Term{subobject}{子对象}，则
\[
 (\sup_i X_i)\cap Y=\sup_i(X_i\cap Y).
\]
\noindent ii) 小的\Term{pseudo-filtered-inductive-limit}{伪滤过归纳极限}（第 I 讲 2.7）\Term{representable}{可表}，并与\Term{finite}{有限}\Term{projective-limit}{逆极限}交换。

此外，若上述条件成立，则小的\Term{filtered-inductive-limit}{滤过归纳极限}是\Term{universal}{普适}的（第 I 讲 2.6）。
\end{SGAresult}
\begin{proof}
$(\mathrm{ii})\Rightarrow(\mathrm{i})$ 显然。为证 $\mathrm{i})\Rightarrow\mathrm{ii})$ 以及补充断言，只须使用下述事实：$\mathcal A$ 是某个适当的环上的\Term{module-category}{模范畴} $\mathcal M$ 的\Term{full-subcategory}{满子范畴}，且\Term{inclusion-functor}{包含函子} $u:\mathcal A\to\mathcal M$ 有\Term{exact}{正合}的\Term{left-adjoint}{左伴随} $v$ [5]。此时验证是平凡的。
\end{proof}
\textbf{0.1.1。}已知 [11]，具有\Term{generator}{生成元}且满足公理 AB 5) 的\Term{abelian-category}{阿贝尔范畴} $\mathcal A$ 有足够多的\Term{injective-object}{内射对象}，即每个\Term{object}{对象}均可\Term{embedding}{嵌入}一个\Term{injective-object}{内射对象}。此外，由已引用的 [5] 的结果，$\mathcal A$ 中的小积\Term{representable}{可表}（公理 AB 3)$^*$）。
\SourcePage{9}
\begin{SGAresult}{命题}{0.2}
设 $\mathcal A$、$\mathcal B$ 是两个\Term{abelian-category}{阿贝尔范畴}，且 $u:\mathcal A\to\mathcal B$ 与 $v:\mathcal B\to\mathcal A$ 是一对\Term{adjoint-functor}{伴随函子}（$u$ 是 $v$ 的\Term{left-adjoint}{左伴随}）。考虑下列性质：

\noindent i) \Term{functor}{函子} $u$ \Term{exact}{正合}。

\noindent ii) \Term{functor}{函子} $v$ 把 $\mathcal B$ 的\Term{injective-object}{内射对象}变为 $\mathcal A$ 的\Term{injective-object}{内射对象}。

1) 总有 $(\mathrm{i})\Rightarrow(\mathrm{ii})$。若还假设 $\mathcal B$ 的每个非零\Term{object}{对象}都有一个到某\Term{injective-object}{内射对象}的非零\Term{morphism}{态射}，则 $(\mathrm{ii})\Leftrightarrow(\mathrm{i})$。

2) 假设：a) $\mathcal B$ 有足够多的\Term{injective-object}{内射对象}；b) 上述两个\Term{equivalence}{等价}条件 (i)、(ii) 之一成立；c) \Term{functor}{函子} $u$ 忠实。则 $\mathcal A$ 有足够多的\Term{injective-object}{内射对象}。
\end{SGAresult}
\begin{proof}证明留给读者。\end{proof}
\textbf{注记 0.2.1。}\Term{abelian-group}{交换群}\Term{category}{范畴}有足够多的\Term{injective-object}{内射对象}。应用引理 0.2，得出含幺环上幺模的任何\Term{category}{范畴}都有足够多的\Term{injective-object}{内射对象}。再应用 [5] 的结果（在 0.1 的证明中已使用）及 0.2，可得任何具有\Term{generator}{生成元}和\Term{exact}{正合}的小\Term{filtered-inductive-limit}{滤过归纳极限}的\Term{abelian-category}{阿贝尔范畴}都有足够多的\Term{injective-object}{内射对象}；这给出了该事实的另一个证明。
\begin{SGAresult}{命题}{0.3}
设 $\mathcal A,\mathcal B,\mathcal C$ 是三个\Term{abelian-category}{阿贝尔范畴}，$u:\mathcal A\to\mathcal B$、$v:\mathcal B\to\mathcal C$ 是两个\Term{additive}{加性}\Term{left-exact-functor}{左正合函子}。假设 $\mathcal A$ 和 $\mathcal B$ 有足够多的\Term{injective-object}{内射对象}。下列性质\Term{equivalence}{等价}：

\noindent i) 存在一个\Term{spectral-functor}{谱函子}，其 $E_2^{p,q}$ 项为 $R^pvR^qu$，且收敛于带有适当\Term{filtered-category}{滤过的} $R^{p+q}vu$。

\noindent ii) \Term{functor}{函子} $u$ 把 $\mathcal A$ 的\Term{injective-object}{内射对象}变为对\Term{functor}{函子} $v$ \Term{acyclic-object}{无环的对象}。
\end{SGAresult}
\begin{proof}
$(\mathrm{i})\Rightarrow(\mathrm{ii})$ 是平凡的，只须把\Term{spectral-functor}{谱函子}应用于\Term{injective-object}{内射对象}。$(\mathrm{ii})\Rightarrow(\mathrm{i})$ 见 [11]。
\end{proof}
\begin{SGAresult}{命题}{0.4}
设 $\mathcal A$、$\mathcal B$ 是两个\Term{abelian-category}{阿贝尔范畴}，$u:\mathcal A\to\mathcal B$ 是\Term{additive}{加性}\Term{left-exact-functor}{左正合函子}。设 $\mathcal M$ 是 $\mathcal A$ 的\Term{object}{对象}集的一个\Term{subset}{子集}，具有下列性质：

1) $\mathcal A$ 的每个\Term{object}{对象}都\Term{embedding}{嵌入} $\mathcal M$ 的某个元素。

2) 若 $X\oplus Y$ 属于 $\mathcal M$，则\Term{object}{对象} $X$ 属于 $\mathcal M$。
\SourcePage{10}
3) 若 $0\to X'\to X\to X''\to0$ 是\Term{exact-sequence}{正合列}，且 $X'$、$X$ 属于 $\mathcal M$，则 $X''$ 属于 $\mathcal M$，且
\[
 0\longrightarrow u(X')\longrightarrow u(X)\longrightarrow u(X'')\longrightarrow0
\]
\Term{exact}{正合}。零\Term{object}{对象}属于 $\mathcal M$。

则每个\Term{injective-object}{内射对象}都属于 $\mathcal M$，且 $\mathcal M$ 的\Term{object}{对象}对 $u$ \Term{acyclic-object}{无环}，即对每个 $q\ne0$ 和 $\mathcal M$ 的每个\Term{object}{对象} $X$，有 $R^qu(X)=0$。（特别地，由 $\mathcal M$ 的\Term{object}{对象}组成的\Term{resolution}{消解}可用来计算 $u$ 的\Term{derived-functor}{导出函子}。）
\end{SGAresult}
证明见 [11] 3.3.1。
\begin{SGAresult}{命题}{0.5}
设 $\mathcal U\subset\mathcal V$ 是两个\Term{universe}{宇宙}，$\mathcal A$（相应地 $\mathcal B$）是一个阿贝尔 $\mathcal U$-\Term{category}{范畴}（相应地 $\mathcal V$-\Term{category}{范畴}），满足相对于 $\mathcal U$（相应地 $\mathcal V$）的公理 AB 5)，且有一个 $\mathcal U$-小（相应地 $\mathcal V$-小）的\Term{generating-family}{生成族}。设 $\varepsilon:\mathcal A\to\mathcal B$ 是\Term{exact}{正合}\Term{fully-faithful}{全忠实}\Term{functor}{函子}。下列条件\Term{equivalence}{等价}：

1) $\mathcal A$ 有一个\Term{generating-family}{生成族} $(X_i)_{i\in I}$，使 $(\varepsilon(X_i))_{i\in I}$ 在 $\mathcal B$ 中也是\Term{generating-family}{生成族}。

$1')$ $\mathcal B$ 的每个\Term{object}{对象}都\Term{isomorphism}{同构}于形如 $\bigoplus_{\alpha\in A}\varepsilon(Y_\alpha)$ 的\Term{object}{对象}的一个商，其中 $A\in\mathcal V$。

在这些条件下，成立：

2) $\varepsilon$ 把 $\mathcal U$-小积变为积（因而与 $\mathcal U$-小\Term{projective-limit}{逆极限}交换）。

此外，在\Term{equivalence}{等价}条件 1) 或 $1')$ 下，下列条件\Term{equivalence}{等价}：

a) 对 $\mathcal A$ 的每个\Term{object}{对象} $Y$，$\varepsilon(Y)$ 的每个\Term{subobject}{子对象}都\Term{isomorphism}{同构}于 $Y$ 的某个\Term{subobject}{子对象}在 $\varepsilon$ 下的像。

$a')$ $\mathcal A$ 有一个\Term{generating-family}{生成族} $(X_i)_{i\in I}$，使 $(\varepsilon(X_i))_{i\in I}$ 在 $\mathcal B$ 中是\Term{generating-family}{生成族}，并且对每个 $i\in I$，$\varepsilon(X_i)$ 的每个\Term{subobject}{子对象}都\Term{isomorphism}{同构}于 $X_i$ 的某个\Term{subobject}{子对象}在 $\varepsilon$ 下的像。

b) $\mathcal B$ 的每个\Term{object}{对象}都\Term{isomorphism}{同构}于形如 $\prod_{\alpha\in A}\varepsilon(Y_\alpha)$ 的\Term{object}{对象}的一个\Term{subobject}{子对象}，其中 $A\in\mathcal V$。

c) $\varepsilon$ 与 $\mathcal U$-小\Term{direct-sum}{直和}交换（因而与 $\mathcal U$-小\Term{inductive-limit}{归纳极限}交换）。

此外，在条件 1) 和 $a')$ 下，有：

d) $\varepsilon$ 把\Term{injective-object}{内射对象}变为\Term{injective-object}{内射对象}。
\end{SGAresult}
\textbf{注记 0.5.1。}在 0.5 中，若 $\mathcal U=\mathcal V$，则条件 1) 和 $a')$ 蕴涵 $\varepsilon$ 是\Term{equivalence-of-categories}{范畴等价}（因为此时 b)、2)、a) 均成立）。
\SourcePage{11}
\textbf{0.5.2。}我们只给出证明的提示。$1)\Leftrightarrow1')$、$1)\Rightarrow2)$ 留给读者。$a)\Rightarrow a')$ 显然。证明 $a')$ 蕴涵 d)。设 $M$ 是 $\mathcal A$ 的一个\Term{injective-object}{内射对象}。对每个 $i\in I$ 和 $X_i$ 的每个\Term{subobject}{子对象} $Y\hookrightarrow X_i$，\Term{homomorphism}{同态} $\operatorname{Hom}(X_i,M)\to\operatorname{Hom}(Y,M)$ 是\Term{surjectivity}{满射}。因此，由 $a')$ 及 $\varepsilon$ 的\Term{fully-faithful}{全忠实性}，对每个\Term{subobject}{子对象} $U\hookrightarrow\varepsilon(X_i)$，\Term{homomorphism}{同态}
\[
 \operatorname{Hom}(\varepsilon(X_i),\varepsilon(M))\longrightarrow
 \operatorname{Hom}(U,\varepsilon(M))
\]
是\Term{surjectivity}{满射}。由于 $\varepsilon(X_i)$ 构成 $\mathcal B$ 的一个\Term{generating-family}{生成族}，$\varepsilon(M)$ 是\Term{injective-object}{内射的} [11]。证明 $a')\Rightarrow b)$。扩大 $X_i$ 族后，可假设每个 $X_i$ 的任何商都\Term{isomorphism}{同构}于某个适当的 $X_j$。对每个 $i\in I$，取一个到\Term{injective-object}{内射对象}的\Term{monomorphism}{单态射} $X_i\hookrightarrow M_i$。族 $(\varepsilon(X_i))_{i\in I}$ 对取商\Term{stable}{稳定}，且由 d)，$\varepsilon(X_i)\hookrightarrow\varepsilon(M_i)$ 是一个到\Term{injective-object}{内射对象}的\Term{monomorphism}{单态射}。由于 $\varepsilon(X_i)$ 是\Term{generating-family}{生成族}，立即可验证 $(\varepsilon(M_i))_{i\in I}$ 是\Term{cogenerating-family}{余生成族}，从而有 b)。

证明 $b)\Rightarrow c)$。设 $(Z_\alpha)_{\alpha\in A}$ 是 $\mathcal A$ 的一族 $\mathcal U$-小\Term{object}{对象}。证明\Term{canonical}{典范}\Term{morphism}{态射}
\[
 \bigoplus_{\alpha\in A}\varepsilon(Z_\alpha)\longrightarrow
 \varepsilon\left(\bigoplus_{\alpha\in A}Z_\alpha\right)
\]
是\Term{isomorphism}{同构}。由于\Term{object}{对象} $\varepsilon(Y)$ 构成\Term{cogenerating-family}{余生成族}，只须证明对 $\mathcal A$ 的每个\Term{object}{对象} $Y$，\Term{homomorphism}{同态}
\[
 \operatorname{Hom}\left(\varepsilon(\bigoplus_a Z_a),\varepsilon(Y)\right)
 \longrightarrow\operatorname{Hom}\left(\bigoplus_a\varepsilon(Z_a),\varepsilon(Y)\right)
\]
是\Term{isomorphism}{同构}；这由 $\varepsilon$ 的\Term{fully-faithful}{全忠实性}得到。最后证明 $c)\Rightarrow a)$。设 $Z$ 是 $\varepsilon(Y)$ 的一个\Term{subobject}{子对象}。由 $1')$，存在 $\mathcal A$ 的一族 $\mathcal U$-小\Term{object}{对象} $Y_\alpha$ 以及一个\Term{epimorphism}{满态射} $\bigoplus_\alpha\varepsilon(Y_\alpha)\to Z$。由于 $\varepsilon$ 与 $\mathcal U$-小\Term{direct-sum}{直和}交换，存在\Term{epimorphism}{满态射} $\varepsilon(Y')\to Z$。记复合\Term{morphism}{态射}为 $u:\varepsilon(Y')\to Z\to\varepsilon(Y)$。有 $Z=\operatorname{Im}(u)$。由 $\varepsilon$ \Term{fully-faithful}{全忠实}，有 $u=\varepsilon(v)$，从而
$Z=\operatorname{Im}(\varepsilon(v))=\varepsilon(\operatorname{Im}(v))$。

\begin{SGAresult}{练习}{0.5.2}
设 $\operatorname{Sex}_{\mathcal V}(\mathcal A)$ 是从 $\mathcal A$ 到 $\mathcal V$-\Term{abelian-group}{交换群}\Term{category}{范畴}、且与 $\mathcal U$-小\Term{inductive-limit}{归纳极限}交换的\Term{contravariant-functor}{反变函子}\Term{category}{范畴}。证明在条件 1) 和 $a')$ 下，从 $\mathcal B$ 到 $\operatorname{Sex}_{\mathcal V}(\mathcal A)$ 的\Term{canonical-functor}{典范函子}是\Term{equivalence-of-categories}{等价}。
\end{SGAresult}

\SGAsection{1}{平坦模}
\begin{SGAresult}{定义}{1.1}
设 $(E,A)$ 是\Term{ringed-topos}{赋环拓扑斯}（第 IV 讲 11.1.1）。若从左（相应地右）$A$-\Term{module-category}{模范畴}到 $E$ 的阿贝尔\Term{sheaf}{层}\Term{category}{范畴}的\Term{functor}{函子} $M\otimes_A-$（相应地 $-\otimes_AM$）\Term{exact}{正合}，则称右（相应地左）$A$-模 $M$ 为\Term{flat-module}{平坦的}。
\end{SGAresult}
\begin{SGAresult}{命题}{1.2}
设 $M$ 是一个 $B$-$A$ \Term{bimodule}{双模}。

1) 下列性质\Term{equivalence}{等价}：
\SourcePage{12}
\noindent i) 模 $M$ 左 $A$-\Term{flat-module}{平坦}；

\noindent ii) 对每个\Term{injective-object}{内射} $B$-模 $I$，左 $A$-模 $\mathscr{H}om_B(M,I)$ 是\Term{injective-object}{内射的}。

2) \Term{flat-module}{平坦模}的\Term{pseudo-filtered-inductive-limit}{伪滤过归纳极限}（第 I 讲 2.7.1）$M$ 是\Term{flat-module}{平坦的}。

3) 最后，若 $M_\bullet=\cdots\to M_{i+1}\to M_i\to\cdots$ 是\Term{flat-module}{平坦模}的\Term{acyclic-complex}{无环复形}（$i<i_0$ 时 $M_i=0$），则对每个模 $F$，\Term{complex}{复形}
\[
 M_\bullet\otimes_AF=\cdots\to M_{i+1}\otimes_AF\longrightarrow M_i\otimes_AF\to\cdots
\]
是\Term{acyclic-complex}{无环的}。
\end{SGAresult}
\begin{proof}
由第 IV 讲 12.12，有\Term{adjunction-isomorphism}{伴随同构}
\[
 \operatorname{Hom}_B(M\otimes_A-,-)\xrightarrow{\sim}
 \operatorname{Hom}_A(-,\mathscr{H}om_B(M,-)).
\]
应用 0.2 即得 $(\mathrm{i})\Leftrightarrow(\mathrm{ii})$。此\Term{adjunction-isomorphism}{伴随同构}也表明\Term{tensor-product}{张量积}与\Term{inductive-limit}{归纳极限}交换。\Term{pseudo-filtered-inductive-limit}{伪滤过归纳极限}的\Term{exactness}{正合性}（0.1）给出第二个断言。要证明 $M_\bullet\otimes_AF$ 是\Term{acyclic-complex}{无环的}，只须证明对每个\Term{injective-object}{内射}阿贝尔\Term{sheaf}{层} $I$，\Term{complex}{复形} $\operatorname{Hom}_{\mathbf Z}^{\bullet}(M_\bullet\otimes_AF,I)$ \Term{acyclic-complex}{无环}。由\Term{adjunction}{伴随}公式，此\Term{complex}{复形}\Term{isomorphism}{同构}于
$\operatorname{Hom}_A^{\bullet}(F,\mathscr{H}om_{\mathbf Z}(M_\bullet,I))$。根据 $(\mathrm{i})\Leftrightarrow(\mathrm{ii})$，\Term{sheaf}{层}\Term{complex}{复形} $\mathscr{H}om_{\mathbf Z}(M_\bullet,I)$ 是由\Term{injective-object}{内射对象}组成的\Term{acyclic-complex}{无环复形}，从而得到结论。
\end{proof}
\begin{SGAresult}{命题}{1.3.1}
设 $(E,A)$ 是\Term{ringed-topos}{赋环拓扑斯}，$H$ 是 $E$ 的\Term{object}{对象}，$M$ 是\Term{flat-module}{平坦} $A/H$-模。则 $j_{H!}M$ 是\Term{flat-module}{平坦} $A$-模。特别地，$A_H$ 左右均\Term{flat-module}{平坦}。
\end{SGAresult}
为明确起见，假设 $M$ 是右 $A/H$-模。对每个左 $A$-模 $P$，有\Term{canonical-isomorphism}{典范同构}（第 IV 讲 12）
\[
 P\otimes_Aj_{H!}M\simeq j_{H!}(P\otimes_{A/H}M).
\]
\Term{functor}{函子} $j_{H!}$ 和 $j_H^*$ \Term{exact}{正合}（第 IV 讲 11.3.1 及 11.12.2），而按假设 $-\otimes_{A/H}M$ \Term{exact}{正合}。因此 $P\mapsto P\otimes_Aj_{H!}M$ \Term{exact}{正合}，$j_{H!}M$ \Term{flat-module}{平坦}。
\begin{SGAresult}{命题}{1.3.2}
（\Term{closed-immersion}{闭浸入}的\Term{projection-formula}{投影公式}。）设 $(E,A)$ 是\Term{ringed-topos}{赋环拓扑斯}，$i:F\to E$ 是 $E$ 的\Term{closed-subtopos}{闭子拓扑斯}。置 $i^*A=A/F$。对每个（右）$A/F$-模 $M$ 和每个（左）$A$-模 $P$，有\Term{functor}{函子}\Term{isomorphism}{同构}
\[
 i_*(M\otimes_{A/F}i^*P)\simeq(i_*M)\otimes_AP.
\]
\end{SGAresult}
\SourcePage{13}
设 $U$ 是 $F$ 的补开子\Term{topos}{拓扑斯}，$j:U\to E$ 是\Term{canonical}{典范}\Term{immersion}{浸入}\Term{morphism}{态射}。有
$j^*(i_*M\otimes_AP)\simeq0\otimes_{A/U}j^*P$（第 IV 讲 12），故 $i_*M\otimes_AP$ 的\Term{support}{支集}包含于 $F$。因此\Term{adjunction-morphism}{伴随态射}
\[
 i_*M\otimes_AP\longrightarrow i_*i^*(i_*M\otimes_AP)
\]
是\Term{isomorphism}{同构}（第 IV 讲 14）。又有 $i^*(i_*M\otimes_AP)\simeq i^*i_*M\otimes_{A/F}i^*P$（第 IV 讲 12）；由于 $i:F\to E$ 是\Term{closed-immersion}{闭浸入}，$i^*i_*M\simeq M$（第 IV 讲 14），从而得到所述\Term{isomorphism}{同构}。
\begin{SGAresult}{推论}{1.3.3}
对每个\Term{flat-module}{平坦} $A/F$-模 $M$，$i_*M$ \Term{flat-module}{平坦}。
\end{SGAresult}
由 1.3.2 和第 IV 讲 14，\Term{functor}{函子} $P\mapsto(i_*M)\otimes_AP$ \Term{exact}{正合}。

\textbf{1.4。}设 $(E,A)$ 是\Term{ringed-topos}{赋环拓扑斯}，$x:P\to E$ 是 $E$ 的一个点（第 IV 讲 6.1），$\operatorname{vois}(x)$ 是 $x$ 的\Term{neighborhood}{邻域}\Term{category}{范畴}（第 IV 讲 6.8）。$\operatorname{vois}(x)$ 的每个\Term{object}{对象} $V$ 对应于 $E$ 的一个\Term{object}{对象}，仍记为 $V$，以及 $E/V$ 的一个点 $x_V:P\to E/V$。此外，$\operatorname{vois}(x)$ 的每个\Term{morphism}{态射} $u:V\to W$ 对应于一个\Term{topos}{拓扑斯}\Term{morphism}{态射}的本质\Term{commutative-diagram}{交换图}（第 IV 讲 6.7）
\[
\begin{tikzcd}
 P\arrow[r,"x_V"]\arrow[dr,"x_W"'] & E/V\arrow[d,"j_u"]\\
 &E/W.
\end{tikzcd}\tag{1.4.1}
\]
设 $N$ 是 $A_x$-模（相应地，\Term{set}{集合}）。由图 (1.4.1) 得到一个 $A_x$-模（相应地，\Term{set}{集合}）的图：
\[
\begin{tikzcd}
 x_W^*x_{W*}N\arrow[r,"\operatorname{ad}_W"]\arrow[d,"\phi(u)"']&N\\
 x_V^*x_{V*}N\arrow[ur,"\operatorname{ad}_V"']&
\end{tikzcd}\tag{1.4.2}
\]
其中 $\operatorname{ad}_W$ 和 $\operatorname{ad}_V$ 是\Term{adjunction-morphism}{伴随态射}。此外，立即可验证 $\phi(uv)=\phi(v)\phi(u)$。因此得到从\Term{filtered-category}{滤过范畴} $\operatorname{vois}(x)^\circ$ 到 $A_x$-\Term{module-category}{模范畴}（相应地，\Term{set}{集合}\Term{category}{范畴}）的\Term{functor}{函子}，以及\Term{homomorphism}{同态}
\[
 \Lambda(N):\varinjlim_{V\in\operatorname{vois}(x)^\circ}x_V^*x_{V*}N\longrightarrow N.
 \tag{1.4.3}
\]
\begin{SGAresult}{命题}{1.5}
对每个 $A_x$-模（相应地，\Term{set}{集合}）$N$，$\Lambda(N)$ 是\Term{isomorphism}{同构}。
\end{SGAresult}
\SourcePage{14}
这一命题是 Deligne 的一个命题（8.2.6）的特例。这里给出直接证明。取底层\Term{set}{集合}后，只须在 $N$ 为\Term{set}{集合}时证明本命题（第 I 讲 2.8）。$\operatorname{vois}(x)$ 的\Term{arrow}{箭头}的\Term{cofiltered-category}{余滤过范畴} $\operatorname{Fl}(\operatorname{vois}(x))$，通过把\Term{arrow}{箭头}映到其靶的\Term{functor}{函子}
$p:\operatorname{Fl}(\operatorname{vois}(x))\to\operatorname{vois}(x)$，\Term{kan-fibration}{纤维化}于 $\operatorname{vois}(x)$ 之上。设 $D$ 是\Term{category}{范畴}，$F:\operatorname{Fl}(\operatorname{vois}(x))\to D$ 是一个\Term{pro-object}{pro-对象}（第 I 讲 8.10）。对每个\Term{neighborhood}{邻域} $V$，记 $F_V$ 为把 $F$ \Term{restriction}{限制}于\Term{fibre-category}{纤维范畴} $\operatorname{vois}(x)/V$ 所得的\Term{pro-object}{pro-对象}。对每个\Term{morphism}{态射} $m:U\to V$，沿 $m$ 的\Term{base-change}{基变换}\Term{functor}{函子}给出从 $F_U$ 到 $F_V$ 的\Term{morphism}{态射}；\Term{canonical}{典范}\Term{pro-object}{pro-对象}\Term{morphism}{态射} $F\to F_V$ 确定一个\Term{pro-object}{pro-对象}\Term{morphism}{态射}
\[
 F\longrightarrow\varprojlim_{V\in\operatorname{vois}(x)^\circ}F_V,\tag{1.5.1}
\]
立即可验证它是\Term{isomorphism}{同构}。把这一观察用于带标点\Term{set}{集合}的\Term{pro-object}{pro-对象}
\[
 (U,V,m)\longmapsto x_V^*(m)=F(U,V,m),\tag{1.5.2}
\]
得到\Term{pro-object}{pro-对象}\Term{isomorphism}{同构}
\[
 F\xrightarrow{\sim}\varprojlim_{\operatorname{vois}(x)}
 \varprojlim_{\operatorname{vois}(x)/V}x_V^*(m).\tag{1.5.3}
\]
因此对每个\Term{set}{集合} $N$，有\Term{bijection}{双射}
\[
\begin{split}
 &\varinjlim_{\operatorname{Fl}(\operatorname{vois}(x))^\circ}
 \operatorname{Hom}(x_V^*(m),N)\\
 &\qquad\simeq\varinjlim_{\operatorname{vois}(x)^\circ}
 \varinjlim_{(\operatorname{vois}(x)/V)^\circ}
 \operatorname{Hom}(x_V^*(m),N)^\circ.
\end{split}\tag{1.5.4}
\]
然而，对固定的 $V$，$\varinjlim_{(\operatorname{vois}(x)/V)^\circ}\operatorname{Hom}(x_V^*(m),N)$ 可等同于 $x_V^*x_{V*}N$。事实上，由第 IV 讲 6.8.1，
\[
 x_V^*x_{V*}N\simeq\varinjlim_{\operatorname{vois}(x_V)^\circ}
 \operatorname{Hom}_{E/V}(W,x_{V*}N),
\]
由\Term{adjunction}{伴随}性得到
\[
 x_V^*x_{V*}N\simeq\varinjlim_{\operatorname{vois}(x_V)^\circ}
 \operatorname{Hom}(x_V^*(w),N),
\]
而 $\operatorname{vois}(x_V)$ \Term{equivalence}{等价}于 $\operatorname{vois}(x)/V$（第 IV 讲 6.7.2）。最后，读者可验证，(1.5.4) 中的\Term{canonical}{典范}\Term{transition-map}{转移映射}
\[
 \varinjlim_{(\operatorname{vois}(x)/U)^\circ}\operatorname{Hom}(x_U^*(m),N)
 \longrightarrow\varinjlim_{(\operatorname{vois}(x)/V)^\circ}\operatorname{Hom}(x_V^*(n),N)
\]
等同于 (1.4.3) 的\Term{canonical}{典范}映射 $x_U^*x_{U*}N\to x_V^*x_{V*}N$。故有\Term{bijection}{双射}
\[
 \varinjlim_{\operatorname{Fl}(\operatorname{vois}(x))^\circ}\operatorname{Hom}(x_V^*(m),N)
 \simeq\varinjlim_{\operatorname{vois}(x)^\circ}x_V^*x_{V*}N.\tag{1.5.5}
\]
映射 $\Lambda(N):\varinjlim_{\operatorname{vois}(x)^\circ}x_V^*x_{V*}N\to N$（1.4.3）与\Term{bijection}{双射} (1.5.5) 的复合，来自如下映射
$\operatorname{Hom}(x_V^*(m),N)\to N$：它把 $r:x_V^*(m)\to N$ 映到 $x_V^*(m)$ 的标点在 $r$ 下的像。要证明命题，只须注意 $\operatorname{Fl}(\operatorname{vois}(x))$ 的每个\Term{object}{对象} $(U,V,m)$ 都有下界 $(U,U,\operatorname{id}_U)$，而 $x_U^*(\operatorname{id}_U)$ 只有一个元素；换言之，从单元素常\Term{pro-object}{pro-对象}到 $F$ 的\Term{canonical}{典范}\Term{morphism}{态射}是\Term{pro-object}{pro-对象}的\Term{isomorphism}{同构}。

\begin{SGAresult}{命题}{1.6}
设 $(E,A)$ 是\Term{ringed-topos}{赋环拓扑斯}，$M$ 是 $A$-模。
\SourcePage{15}
1) 若 $M$ \Term{flat-module}{平坦}，则对 $E$ 的每个点 $x:P\to E$，$A_x$-模 $M_x$ \Term{flat-module}{平坦}。

2) 设 $(x_i)_{i\in I}$ 是 $E$ 的一个\Term{conservative-family}{保守点族}，且对每个 $i\in I$，$M_{x_i}$ 是\Term{flat-module}{平坦} $A_{x_i}$-模。则 $M$ \Term{flat-module}{平坦}。
\end{SGAresult}
\begin{SGAresult}{引理}{1.6.1}
设 $M$ 是\Term{flat-module}{平坦模}，$V$ 是 $E$ 的\Term{object}{对象}。$A/V$-模 $j_V^*M$ \Term{flat-module}{平坦}。
\end{SGAresult}
须证明\Term{functor}{函子} $P\mapsto P\otimes_{A/V}j_V^*M$ \Term{exact}{正合}。由于\Term{extension-by-zero}{零延拓}\Term{functor}{函子} $j_{V!}$ \Term{exact}{正合}且忠实（第 IV 讲 11.3.1），只须证明 $P\mapsto j_{V!}(P\otimes_{A/V}j_V^*M)$ \Term{exact}{正合}。有\Term{functor}{函子}\Term{isomorphism}{同构}
\[
 j_{V!}(P\otimes_{A/V}j_V^*M)\simeq j_{V!}(P)\otimes_AM
\]
（第 IV 讲 12.11），而\Term{functor}{函子} $P\mapsto j_{V!}(P)\otimes_AM$ \Term{exact}{正合}（第 IV 讲 11.3.1）。

\textbf{1.6.2。}证明 1)。为明确起见，假设 $M$ 是\Term{flat-module}{平坦}左 $A$-模，设 $x:P\to E$ 是 $E$ 的点。证明从右 $A_x$-\Term{module-category}{模范畴}到\Term{abelian-group}{交换群}\Term{category}{范畴}的\Term{functor}{函子} $N\mapsto N\otimes_{A_x}M_x$ \Term{exact}{正合}；只须证明它\Term{left-exact}{左正合}。沿用 1.4 的记号，设 $V$ 是 $x$ 的\Term{neighborhood}{邻域}，$j_V:E/V\to E$ 是\Term{localization}{局部化}\Term{morphism}{态射}。对每个 $A_x$-模 $N$，有
\[
\begin{split}
 x_V^*x_{V*}N\otimes_{A_x}M_x
 &\simeq x_V^*x_{V*}N\otimes_{A_x}x_V^*(j_V^*M)\\
 &\simeq x_V^*(x_{V*}N\otimes_{A_x}j_V^*M)
\end{split}
\]
（第 IV 讲 12.11）。故 $N\mapsto x_V^*x_{V*}N\otimes_{A_x}M_x$ \Term{left-exact}{左正合}。于是，由 1.5，\Term{functor}{函子} $N\mapsto N\otimes_{A_x}M_x$ 是\Term{left-exact-functor}{左正合函子}的\Term{filtered-inductive-limit}{滤过归纳极限}，因而\Term{left-exact}{左正合}。证明 2)。设 $0\to P'\to P\to P''\to0$ 是 $A$-模的\Term{exact-sequence}{正合列}，须证明
\[
 0\longrightarrow P'\otimes_AM\longrightarrow P\otimes_AM
 \longrightarrow P''\otimes_AM\longrightarrow0
\]
\Term{exact}{正合}。只须对每个 $i\in I$，证明取 $x_i$ 处\Term{fiber}{纤维}所得的列\Term{exact}{正合}（第 IV 讲 6）。该\Term{fiber}{纤维}列是（第 IV 讲 13.5）
\[
 0\to P'_{x_i}\otimes_{A_{x_i}}M_{x_i}\to P_{x_i}\otimes_{A_{x_i}}M_{x_i}
 \to P''_{x_i}\otimes_{A_{x_i}}M_{x_i}\to0;
\]
它因 $M_{x_i}$ \Term{flat-module}{平坦}而\Term{exact}{正合}。
\begin{SGAresult}{命题}{1.7}
设 $(E,A)$ 是\Term{ringed-topos}{赋环拓扑斯}，$u:E'\to E$ 是\Term{topos}{拓扑斯}\Term{morphism}{态射}，置 $A'=u^*A$。设 $M$ 是右（相应地左）$A'$-模。下列条件\Term{equivalence}{等价}：

\noindent i) 从左（相应地右）$A$-\Term{module-category}{模范畴}到 $E'$ 的阿贝尔\Term{sheaf}{层}\Term{category}{范畴}的\Term{functor}{函子} $N\mapsto M\otimes_{A'}u^*N$（相应地 $N\mapsto u^*N\otimes_{A'}M$）\Term{exact}{正合}。

\noindent ii) $M$ 是\Term{flat-module}{平坦} $A'$-模。
\end{SGAresult}
$\mathrm{ii})\Rightarrow\mathrm{i})$ 显然。证明 $\mathrm{i})\Rightarrow\mathrm{ii})$。这里只在 $E'$ 有一个\Term{conservative-family}{保守点族} $(x_i)_{i\in I}$ 的情形下证明；一般情形在 8.2.7 中处理。这一特例涵盖大部分应用；利用 1.6 和第 IV 讲 11.3.1，可归结为 $E'$ 是点\Term{topos}{拓扑斯}的情形，于是\Term{morphism}{态射} $u$（记为 $x$）是 $E$ 的点。只讨论 $M$ 是右 $A'$-模的情形。
\SourcePage{16}
左 $A'$-模的情形可通过取反环归结于此。设 $V$ 是 $x$ 的\Term{neighborhood}{邻域}，$V_x$ 是其在 $x$ 处的\Term{fiber}{纤维}。沿用 1.4 的记号，有\Term{topos}{拓扑斯}\Term{morphism}{态射}的本质\Term{commutative-diagram}{交换图}
\[
\begin{tikzcd}
 P\arrow[r,"j"]\arrow[dr,"x_V"']&P/V_x\arrow[r,"j_{V_x}"]\arrow[d,"x/V"]&P\arrow[d,"x"]\\
 &E/V\arrow[r]&E.
\end{tikzcd}
\]
其中 $x/V$ 是由 $x$ 在 $V$ 上\Term{localization}{局部化}得到的\Term{morphism}{态射}（第 IV 讲 5.10），$j$ 是由 $V_x$ 的标点得到的\Term{morphism}{态射}。证明\Term{functor}{函子} $N'\mapsto M\otimes_{A'}x_V^*N'$ \Term{exact}{正合}。有 $x_V^*=j^*(x/V)^*$ 及 $\operatorname{id}=j^*j_{V_x}^*$，故有\Term{canonical-isomorphism}{典范同构}
\[
 M\otimes_{A'}x_V^*N'\simeq j^*\bigl(j_{V_x}^*M\otimes_{A'/V_x}(x/V)^*N'\bigr).
\]
由于 $j^*$ \Term{exact}{正合}，而 $j_{V_x!}$ \Term{exact}{正合}且忠实（第 IV 讲 11.3.1），只须证明
\[
 N'\longmapsto j_{V_x!}\bigl(j_{V_x}^*M\otimes(x/V)^*N'\bigr)
\]
\Term{exact}{正合}；由第 IV 讲 12.11，只须证明 $N'\mapsto M\otimes_{A'}x^*j_{V!}N'$ \Term{exact}{正合}，而这是由假设得到的。对每个 $A'$-模 $Q$，有\Term{functor}{函子}\Term{isomorphism}{同构}（1.5）
\[
 Q\simeq\varinjlim_{\operatorname{vois}(x)^\circ}x_V^*x_{V*}Q.
\]
由上可知，$Q\mapsto M\otimes_{A'}Q$ 是\Term{left-exact-functor}{左正合函子}的\Term{filtered-inductive-limit}{滤过归纳极限}，因而\Term{left-exact}{左正合}，所以 $M$ \Term{flat-module}{平坦}。
\begin{SGAresult}{推论}{1.7.1}
设 $u:(E',A')\to(E,A)$ 是\Term{ringed-topos}{赋环拓扑斯}\Term{morphism}{态射}，$M$ 是\Term{flat-module}{平坦} $A$-模。则 $u^*M$ 是\Term{flat-module}{平坦} $A'$-模。
\end{SGAresult}
由 1.7 的判据及公式第 IV 讲 13.4.5 得到。
\begin{SGAresult}{推论}{1.7.2}
设 $u:(E',A')\to(E,A)$ 是\Term{ringed-topos}{赋环拓扑斯}\Term{morphism}{态射}。下列条件\Term{equivalence}{等价}：

\noindent i) 右（相应地左）$u^{-1}(A)$-模 $A'$ \Term{flat-module}{平坦}。

\noindent ii) 从左（相应地右）$A$-\Term{module-category}{模范畴}到左（相应地右）$A'$-\Term{module-category}{模范畴}的\Term{functor}{函子} $u^*$ \Term{exact}{正合}。
\end{SGAresult}
\Term{equivalence}{等价}性由 1.7 和 $u^*$ 的定义（第 IV 讲 13.2.1）得到。
\begin{SGAresult}{定义}{1.8}
具有 1.7.2 的\Term{equivalence}{等价}性质的\Term{ringed-topos}{赋环拓扑斯}\Term{morphism}{态射}，称为左\Term{flat-morphism}{平坦}（相应地右\Term{flat-morphism}{平坦}）\Term{topos}{拓扑斯}\Term{morphism}{态射}。
\end{SGAresult}
\begin{SGAresult}{命题}{1.9}
设 $\mathcal U\subset\mathcal V$ 是两个\Term{universe}{宇宙}，$C$ 是 $\mathcal U$-\Term{site}{景}，$\widetilde C_{\mathcal U}$ 与 $\widetilde C_{\mathcal V}$ 分别是 $\mathcal U$-\Term{set}{集合}\Term{sheaf}{层}和 $\mathcal V$-\Term{set}{集合}\Term{sheaf}{层}的\Term{category}{范畴}，$\varepsilon:\widetilde C_{\mathcal U}\to\widetilde C_{\mathcal V}$ 是\Term{canonical}{典范}\Term{inclusion-functor}{包含函子}，$A$ 是 $C$ 上的 $\mathcal U$-环\Term{sheaf}{层}。
\SourcePage{17}
1) $\varepsilon$ 与模的 $\mathcal U$-小\Term{inductive-limit}{归纳极限}和\Term{projective-limit}{逆极限}交换。在\Term{module-category}{模范畴}上，$\varepsilon$ 是\Term{conservative-functor}{保守}且\Term{fully-faithful}{全忠实}的。

2) 对 $\widetilde C_{\mathcal U}$ 的每个\Term{object}{对象} $H$，有\Term{canonical-isomorphism}{典范同构} $\varepsilon(A_H)\simeq\varepsilon(A)_{\varepsilon(H)}$。

3) 每个左或右\Term{injective-object}{内射} $A$-模都被 $\varepsilon$ 变为\Term{injective-object}{内射} $\varepsilon(A)$-模。

4) 每个右或左\Term{flat-module}{平坦} $A$-模都被 $\varepsilon$ 变为\Term{flat-module}{平坦} $\varepsilon(A)$-模。

5) \Term{functor}{函子} $\varepsilon$ 与\Term{tensor-product}{张量积}和\Term{hom-sheaf}{同态层}的构造交换。
\end{SGAresult}
\Term{associated-sheaf}{伴随层}的构造不依赖\Term{universe}{宇宙}（第 II 讲 3.6）。由\Term{sheaf}{层}\Term{category}{范畴}中\Term{inductive-limit}{归纳极限}和\Term{projective-limit}{逆极限}的构造（第 II 讲 4.1、6.4），$\varepsilon$ 与\Term{set}{集合}或模的 $\mathcal U$-小\Term{projective-limit}{逆极限}交换，从而有 1)。证明 2)。取 $\widetilde C_{\mathcal U}$ 的一个 $\mathcal U$-小\Term{generating-subcategory}{生成子范畴}，可归结到 $C$ 是 $\mathcal U$-小的情形（第 III 讲 4.1）。由第 II 讲 6.5，$H$ 生成的\Term{free-module}{自由} $A$-模，可先构造 $H$ 生成的自由 $A$-模\Term{presheaf}{预层}，再取\Term{associated-sheaf}{伴随层}得到；这两个运算均与\Term{universe}{宇宙}扩大交换，后者见第 II 讲 3.6。为证 3)，根据 0.5，只须证明 $\mathcal U$-\Term{sheaf}{层}的每个 $\mathcal V$-子\Term{sheaf}{层}都是 $\mathcal U$-\Term{sheaf}{层}，这是显然的。证明 4)。设 $M$ 是\Term{flat-module}{平坦} $A$-模。只须证明对每个\Term{injective-object}{内射} $\varepsilon(A)$-模 $J$，阿贝尔\Term{sheaf}{层} $\mathscr{H}om_{\varepsilon(A)}(\varepsilon(M),J)$ \Term{injective-object}{内射}（1.2）。根据 0.5，$J$ 是形如 $\prod_\alpha\varepsilon(I_\alpha)$ 的\Term{object}{对象}的\Term{subobject}{子对象}，因而是其\Term{direct-factor}{直因子}，其中 $I_\alpha$ \Term{injective-object}{内射}。所以可假设 $J=\varepsilon(I)$，$I$ 为\Term{injective-object}{内射对象}。若\Term{homomorphism}{同态}
\[
 \varepsilon(\mathscr{H}om(M,I))\longrightarrow
 \mathscr{H}om_{\varepsilon A}(\varepsilon M,\varepsilon I)
\]
是\Term{isomorphism}{同构}，则由 3) 知 $\mathscr{H}om_{\varepsilon(A)}(\varepsilon(M),\varepsilon(I))$ \Term{injective-object}{内射}。因此只剩下证明 5)。\Term{tensor-product}{张量积}的构造与 $\varepsilon$ 交换，由第 IV 讲 12.6 及\Term{associated-sheaf}{伴随层}的构造与\Term{universe}{宇宙}扩大交换（第 II 讲 3.6）得到。对每个\Term{object}{对象} $X\in\widetilde C_{\mathcal U}$ 以及其中每对 $A$-模 $M,N$，由 $\varepsilon$ 的\Term{fully-faithful}{全忠实性}，有
\[
\begin{split}
 \operatorname{Hom}_{\widetilde C_{\mathcal V}}(\varepsilon X,\mathscr{H}om_A(M,N))
 &\simeq\operatorname{Hom}_{\widetilde C_{\mathcal U}}(X,\mathscr{H}om_A(M,N))\\
 &\simeq\operatorname{Hom}_A(\mathbf Z_X\otimes_{\mathbf Z}M,N).
\end{split}
\]
再由 $\varepsilon$ 的\Term{fully-faithful}{全忠实性}及以上结果，
\[
\begin{split}
 \operatorname{Hom}_A(\mathbf Z_X\otimes_{\mathbf Z}M,N)
 &\simeq\operatorname{Hom}_{\varepsilon A}(\varepsilon(\mathbf Z_X\otimes_{\mathbf Z}M),\varepsilon N)\\
 &\simeq\operatorname{Hom}_{\varepsilon A}(\varepsilon\mathbf Z_X\otimes_{\mathbf Z}\varepsilon M,\varepsilon N).
\end{split}
\]
于是利用 2) 及第 IV 讲 12.14，最终得到\Term{isomorphism}{同构}
\[
\begin{split}
 &\operatorname{Hom}_{\widetilde C_{\mathcal V}}(\varepsilon X,\varepsilon\mathscr{H}om_A(M,N))\\
 &\qquad\simeq\operatorname{Hom}_{\widetilde C_{\mathcal V}}(\varepsilon X,\mathscr{H}om_{\varepsilon A}(\varepsilon M,\varepsilon N)),
\end{split}
\]
从而得到所述\Term{isomorphism}{同构}。

\SourcePage{18}
\textbf{1.10。}设 $E$ 是\Term{topos}{拓扑斯}，$\mathcal U=(U_i\to X)_{i\in I}$ 是一族同靶的小\Term{morphism}{态射}。对每个\Term{finite}{有限}集 $\Delta_n=[0,\ldots,n]$，置
\[
 S_n(\mathcal U)=\coprod_{f:\Delta_n\to I}U_{f(1)}\times_X\cdots\times_XU_{f(n)},
\]
其中和遍历从 $\Delta_n$ 到 $I$ 的映射集。设 $m,n$ 是整数，$g:\Delta_m\to\Delta_n$ 是映射。定义\Term{morphism}{态射}
\[
 s(g):S_n(\mathcal U)\longrightarrow S_m(\mathcal U)\tag{1.10.1}
\]
如下：对每个 $f:\Delta_n\to I$，$s(g)$ 在 $U_{f(1)}\times_X\cdots\times_XU_{f(n)}$ 上的\Term{restriction}{限制}，是\Term{canonical}{典范}包含
\[
 U_{f(g(1))}\times_X\cdots\times_XU_{f(g(m))}\hookrightarrow S_m(\mathcal U)
\]
与唯一\Term{morphism}{态射}
\[
 s_f(g):U_{f(1)}\times_X\cdots\times_XU_{f(n)}
 \longrightarrow U_{f(g(1))}\times_X\cdots\times_XU_{f(g(m))}
\]
的复合，其中 $s_f(g)$ 满足：对每个 $i\in\Delta_m$，
\[
 \operatorname{pr}_{f(g(i))}s_f(g)=\operatorname{pr}_{f(g(i))}
\]
（$\operatorname{pr}_a$ 表示第 $a$ 个\Term{projection}{投影}）。于是得到从\Term{finite}{有限}集\Term{category}{范畴}到 $E$ 的\Term{contravariant-functor}{反变函子} $\Delta_n\mapsto S_n$，即 $E$ 中\Term{object}{对象}的\Term{semisimplicial-complex}{半单纯复形} $S_\bullet(\mathcal U)$。此\Term{complex}{复形}带有到 $X$ 的\Term{canonical}{典范}\Term{augmentation}{增广}。从 $E$ 到\Term{category}{范畴} $C$ 的每个\Term{functor}{函子}都把 $S_\bullet(\mathcal U)$ 变为 $C$ 的\Term{simplicial-object}{单纯对象}。特别地，若 $A$ 是 $E$ 的环\Term{object}{对象}，“生成的\Term{free-module}{自由} $A$-模”\Term{functor}{函子}把 $S_\bullet(\mathcal U)$ 变为一个\Term{augmentation}{增广}到 $A_X$ 的 $A$-\Term{bimodule}{双模}的\Term{simplicial-complex}{单纯复形}，记为 $A_\bullet(\mathcal U)$。有
\[
 A_n(\mathcal U)=\bigoplus_{f:\Delta_n\to I}
 A_{U_{f(1)}\times_X\cdots\times_XU_{f(\square)} }.\tag{1.10.2}
\]
\SGAsourcegap{PDF18式1.10.2，最后一个f的括号内缺失指标；两种原本均空白。}
此\Term{complex}{复形}常记为
\[
\begin{tikzcd}[column sep=large]
 \cdots\arrow[r,shift left=2,"s_0"]\arrow[r,"s_1" description]\arrow[r,shift right=2,"s_2"']&
 \displaystyle\bigoplus_{i,j}A_{U_i\times_XU_j}
 \arrow[r,shift left=1,"s_0"]\arrow[r,shift right=1,"s_1"']&
 \displaystyle\bigoplus_i A_{U_i}\arrow[r]&A_X.
\end{tikzcd}\tag{1.10.3}
\]
(1.10.3) 的\Term{arrow}{箭头}（除最后的\Term{augmentation}{增广}\Term{arrow}{箭头}外）表示\Term{simplicial-complex}{单纯复形} $A_\bullet(\mathcal U)$ 的\Term{face-operator}{面算子}，即对应于从 $\Delta_n$ 到 $\Delta_{n+1}$ 的递增\Term{injectivity}{单射}的\Term{morphism}{态射}（$s_i$ 避开整数 $i$）。给 $A_\bullet(\mathcal U)$ 配上一个\Term{augmentation}{增广}到 $A_X$ 的\Term{differential}{微分}\Term{complex}{复形}：
\[
 \cdots\xrightarrow{d}\bigoplus_{i,j}A_{U_i\times_XU_j}
 \xrightarrow{d}\bigoplus_i A_{U_i}\longrightarrow A_X,\tag{1.10.4}
\]
其中置
\[
 d=\sum_i(-1)^is_i.\tag{1.10.5}
\]
\SourcePage{19}
\begin{SGAresult}{命题}{1.11}
若族 $\mathcal U=(U_i\to X)_{i\in I}$ 是\Term{epimorphic-family}{满态射族}，则\Term{differential}{微分}\Term{complex}{复形} (1.10.4) \Term{acyclic-complex}{无环}，并给出 $A_X$ 的一个\Term{resolution}{消解}。
\end{SGAresult}
记 $\mathbf Z$ 为值为 $\mathbf Z$ 的常\Term{sheaf}{层}。由“生成的\Term{free-module}{自由} $A$-模”\Term{functor}{函子}的定义，对 $E$ 的每个\Term{object}{对象} $H$ 有
\[
 A_H\simeq\mathbf Z_H\otimes_{\mathbf Z}A,
 \qquad A_\bullet(\mathcal U)\simeq\mathbf Z_\bullet(\mathcal U)\otimes_{\mathbf Z}A.
\]
由于 $\mathbf Z_\bullet(\mathcal U)$ 的各项是\Term{flat-module}{平坦} $\mathbf Z$-模（1.3.1），只须在 $A=\mathbf Z$ 时证明命题。

先设 $E$ 是\Term{set}{集合}\Term{topos}{拓扑斯}。则\Term{augmentation}{增广}\Term{complex}{复形} $S_\bullet(\mathcal U)$ 是如下类型的\Term{augmentation}{增广}\Term{complex}{复形}的\Term{direct-sum}{直和}：
\[
\begin{tikzcd}[column sep=large]
\cdots\arrow[r,shift left=3]\arrow[r,shift left=1]\arrow[r,shift right=1]\arrow[r,shift right=3]&
 S\times S\times S\arrow[r,shift left=2]\arrow[r]\arrow[r,shift right=2]&
 S\times S\arrow[r,shift left=1]\arrow[r,shift right=1]&S\arrow[r]&\{e\}.
\end{tikzcd}
\]
这里 $\{e\}$ 是单元素集。每个这样的\Term{augmentation}{增广}\Term{complex}{复形}都是\Term{homotopically-trivial}{同伦平凡}的。因此 $S_\bullet(\mathcal U)$ 是\Term{homotopically-trivial}{同伦平凡}的\Term{augmentation}{增广}\Term{complex}{复形}，其\Term{homology}{同调}平凡，故在此情形下命题成立。

现设 $p:\operatorname{Ens}\to E$ 是 $E$ 的一个点。由于 $\mathbf Z_\bullet(\mathcal U)$ 的构造与\Term{topos}{拓扑斯}\Term{morphism}{态射}的\Term{inverse-image}{逆像}\Term{functor}{函子}交换，
$p^*(\mathbf Z_\bullet(\mathcal U))\simeq\mathbf Z_\bullet(p^*(\mathcal U))$ 是 $\mathbf Z_{p^*X}\simeq p^*(\mathbf Z_X)$ 的\Term{resolution}{消解}。所以若 $E$ 有足够多的\Term{fiber-functor}{纤维函子}，命题成立（第 IV 讲 4.6）。特别地，当 $E$ 是\Term{presheaf}{预层}\Term{topos}{拓扑斯} $\widehat C$ 时属于此情形，因为对 $C$ 的每个\Term{object}{对象} $X$，$\Gamma(X,-)$ 都是\Term{fiber-functor}{纤维函子}。一般情形下，$E$ \Term{equivalence}{等价}于一个小\Term{site}{景} $C$ 上的\Term{sheaf}{层}\Term{topos}{拓扑斯}（第 IV 讲 1），且\Term{epimorphic-family}{满态射族} $\mathcal U=(U_i\to X)_{i\in I}$ 是某\Term{epimorphic-family}{满态射族} $\mathcal U'=(U'_i\to X')_{i\in I}$ 在“\Term{associated-sheaf}{伴随层}”\Term{functor}{函子}下的像。因此，$\mathbf Z_\bullet(\mathcal U)\simeq\mathbf Z_\bullet(\mathcal U')$ 是 $\mathbf Z_{X'}\simeq\mathbf Z_X$ 的\Term{resolution}{消解}。

\SGAsection{2}{Čech 上同调；上同调记号}
\textbf{2.1。一般记号。}

\textbf{2.1.1。}设 $(E,A)$ 是\Term{ringed-topos}{赋环拓扑斯}，$M,N$ 是两个 $A$-模（为明确起见，取左模）。记 $\operatorname{Ext}_A^q(E;M,N)$ 为\Term{functor}{函子} $\operatorname{Hom}_A(M,-)$ 的第 $q$ 个右\Term{derived-functor}{导出函子}在 $N$ 上的值 [11]。\Term{functor}{函子} $\operatorname{Ext}_A^q(E;M,N)$，$q\in\mathbf N$，关于变量 $N$ 构成一个 $\delta$-\Term{functor}{函子}，关于 $M$ 也是一个\Term{contravariant}{反变} $\delta$-\Term{functor}{函子}。按定义，$\operatorname{Ext}_A^0(E;M,N)=\operatorname{Hom}_A(M,N)$。
\SourcePage{20}
\textbf{2.1.2。}设 $X$ 是 $E$ 的\Term{object}{对象}。当 $M=A_X$，即 $X$ 生成的\Term{free-module}{自由} $A$-模（第 IV 讲 12）时，置
$\operatorname{Ext}_A^q(E;A_X,N)=H^q(X,N)$。在这一新记号中，环 $A$ 不再出现。这不会造成混淆，因为将证明 $H^q(X,-)$ 的构造与\Term{restriction-of-scalars}{标量限制}交换（3.5）。\Term{functor}{函子} $H^q(X,-)$ 是\Term{functor}{函子} $\operatorname{Hom}_A(A_X,-)=\operatorname{Hom}_E(X,-)$（也记为 $\Gamma(X,-)$）的第 $q$ 个右\Term{derived-functor}{导出函子}。当 $M=A$ 时，置 $\operatorname{Ext}_A^q(E;A,N)=H^q(E,N)$。

\textbf{2.2。\Term{localization}{局部化}。}设 $X$ 是 $E$ 的\Term{object}{对象}，$j_X:E/X\to E$ 是\Term{localization}{局部化}\Term{morphism}{态射}（第 IV 讲 8）。模上的\Term{functor}{函子} $j_X^*$ \Term{exact}{正合}（4.11），并有\Term{exact}{正合}的左\Term{adjoint-functor}{伴随函子} $j_{X!}$（4.11）。因此它把\Term{injective-object}{内射模}变为\Term{injective-object}{内射模}。故对 $E$ 的可变 $A$-模 $N$ 和 $E/X$ 的可变 $A|X$-模 $M$，有\Term{functor}{函子}\Term{isomorphism}{同构}
\[
 \operatorname{Ext}_{A|X}^q(E/X;M,j_X^*N)
 \simeq\operatorname{Ext}_A^q(E;j_{X!}M,N).\tag{2.2.1}
\]
特别地，有\Term{canonical-isomorphism}{典范同构} $H^q(E/X,j_X^*N)\simeq H^q(X,N)$。
对 $E$ 的每个\Term{object}{对象} $X$ 及每对 $A$-模 $M,N$，置
\[
 \operatorname{Ext}_A^q(X;M,N)\simeq
 \operatorname{Ext}_{A|X}^q(E/X;M|X,N|X).\tag{2.2.2}
\]
由上可知，\Term{functor}{函子} $\operatorname{Ext}_A^q(X;M,-)$ 是\Term{functor}{函子} $N\mapsto\operatorname{Hom}_{A|X}(M|X,N|X)$ 的\Term{derived-functor}{导出函子}。\Term{functor}{函子} $(M,N)\mapsto\operatorname{Ext}_A^q(X;M,N)$，$q>0$，关于每个变量都构成 $\delta$-\Term{functor}{函子}。

\textbf{2.3。\Term{presheaf}{预层}\Term{topos}{拓扑斯}的情形；\Term{covering}{覆盖}的\Term{cohomology}{上同调}。}

\textbf{2.3.1。}设 $C$ 是带环\Term{presheaf}{预层} $A$ 的\Term{small-category}{小范畴}，$\widehat C$ 是 $C$ 上的\Term{presheaf}{预层}\Term{topos}{拓扑斯}，$X$ 是 $\widehat C$ 的\Term{representable}{可表}\Term{object}{对象}。把 $A$-模 $M$ 映到群 $\Gamma(X,M)=M(X)$ 的\Term{functor}{函子}\Term{exact}{正合}（第 I 讲 3）。因此对每个 $q>0$ 和每个 $A$-模 $M$，有 $H^q(X,M)=0$；也就是说，$A_X$ 是\Term{projective-module}{投射} $A$-模。

\textbf{2.3.2。}设 $S$ 是 $C$ 上的\Term{presheaf}{预层}。对每个 $A$-模 $M$，有\Term{canonical-isomorphism}{典范同构}（第 I 讲 2）
\[
 \Gamma(S,M)\simeq\varprojlim_{C/S}M|S.
\]
此外，对每个\Term{injective-object}{内射} $A$-模 $M$，$A|S$-模 $M|S$ \Term{injective-object}{内射}（2.2）。因此群 $H^q(S,M)$ 是\Term{functor}{函子} $\varprojlim_{C/S}$ 的右\Term{derived-functor}{导出函子}在 $M|S$ 上的值。记这些\Term{derived-functor}{导出函子}为 $\varprojlim_{C/S}^{q}$，有\Term{canonical-isomorphism}{典范同构}
\[
 H^q(S,M)\simeq\varprojlim_{C/S}^{q}M|S.
\]
特别地，有\Term{canonical-isomorphism}{典范同构}
\[
 H^q(\widehat C,M)\simeq\varprojlim_{C/S}^{q}M.
\]
\SourcePage{21}
\textbf{2.3.3。}设 $X$ 是 $C$ 的\Term{object}{对象}，$\mathcal U=(U_i\to X)_{i\in I}$ 是 $C$ 的一族\Term{morphism}{态射}，且每个 $U_i\to X$ 都\Term{quarrable}{可作基变换}（第 I 讲 10）。记 $A_\bullet$ 为 1.10 中研究的\Term{simplicial-complex}{单纯复形}：
\[
\begin{tikzcd}[column sep=large]
 A_\bullet=\cdots\arrow[r,shift left=2]\arrow[r]\arrow[r,shift right=2]&
 \displaystyle\coprod_{(i,j)\in I\times I}A_{U_i\times_XU_j}
 \arrow[r,shift left=1]\arrow[r,shift right=1]&\displaystyle\coprod_{i\in I}A_{U_i}.
\end{tikzcd}
\]
对每个 $A$-模 $M$，置 $C^\bullet(\mathcal U,M)=\operatorname{Hom}_A(A_\bullet,M)$：
\[
\begin{tikzcd}[column sep=large]
 C^\bullet(\mathcal U,M):\displaystyle\prod_{i\in I}M(U_i)
 \arrow[r,shift left=1]\arrow[r,shift right=1]&
 \displaystyle\prod_{(i,j)\in I\times I}M(U_i\times_XU_j)
 \arrow[r,shift left=2]\arrow[r]\arrow[r,shift right=2]&\cdots.
\end{tikzcd}\tag{2.3.3.1}
\]
置 $H^q(\mathcal U,M)=H^q(C^\bullet(\mathcal U,M))$。
\begin{SGAresult}{命题}{2.3.4}
1) 沿用 2.3.3 的记号，设 $R\hookrightarrow X$ 是族 $(U_i\to X)_{i\in I}$ 生成的\Term{sieve}{筛}。有\Term{canonical-isomorphism}{典范同构}
\[
 H^q(\mathcal U,M)\simeq H^q(R,M).
\]
2) \Term{functor}{函子} $H^q(\mathcal U,-)$ 与\Term{restriction-of-scalars}{标量限制}交换。
\end{SGAresult}
由于 $R$ 是 $\widehat C$ 中 $X$ 的\Term{subobject}{子对象}，\Term{fiber-product}{纤维积}
$U_{i_1}\times_RU_{i_2}\times_R\cdots\times_RU_{i_n}$ 与
$U_{i_1}\times_XU_{i_2}\times_X\cdots\times_XU_{i_n}$ \Term{canonical-isomorphism}{典范同构}。由 1.11，$A_{\mathcal U\bullet}$ 是 $A_R$ 的\Term{resolution}{消解}；由 2.3.1，其各项为\Term{projective-module}{投射模}。故 $C^\bullet(\mathcal U,M)$ 的\Term{cohomology-group}{上同调群}\Term{canonical-isomorphism}{典范同构}于 $\operatorname{Ext}_A^q(\widehat C;A_R,M)$，从而得到所述\Term{isomorphism}{同构}。断言 2) 直接来自 (2.3.3.1) 对 $C^\bullet(\mathcal U,M)$ 的描述。
\begin{SGAresult}{推论}{2.3.5}
设 $\mathcal U=(U_i\to X,i\in I)$ 与 $\mathcal U'=(U'_j\to X,j\in J)$ 是两个靶为 $X$ 的\Term{morphism}{态射}族。设
$\phi=(\phi:I\to J,m_i:U_i\to U'_{\phi(i)})$ 及
$\phi'=(\phi':I\to J,m'_i:U_i\to H'_{\phi(i)})$ 是从 $\mathcal U$ 到 $\mathcal U'$ 的两个（在 $X$ 上方的）\Term{morphism}{态射}。\Term{morphism}{态射} $\phi$、$\phi'$ 诱导\Term{morphism}{态射}
\[
 \phi^q,\phi'^q:H^q(\mathcal U,M)\longrightarrow H^q(\mathcal U',M).
\]
$\phi^q$ 与 $\phi'^q$ 相等。特别地，若族 $\mathcal U$、$\mathcal U'$ \Term{equivalence}{等价}（即存在从 $\mathcal U$ 到 $\mathcal U'$ 及从 $\mathcal U'$ 到 $\mathcal U$ 的\Term{morphism}{态射}），则 $A$-模 $H^q(\mathcal U,M)$ 与 $H^q(\mathcal U',M)$ \Term{canonical-isomorphism}{典范同构}。
\end{SGAresult}
\begin{proof}
事实上，设 $R,R'$ 是 $X$ 中分别由 $\mathcal U,\mathcal U'$ 生成的\Term{sieve}{筛}。\Term{morphism}{态射} $\phi,\phi'$ 定义从 $R$ 到 $R'$ 的同一\Term{morphism}{态射}，因而在\Term{projective-resolution}{投射消解} $A_{\mathcal U\bullet}$ 与 $A_{\mathcal U'\bullet}$ 之间诱导两个\Term{homotopy}{同伦}\Term{morphism}{态射}。
\end{proof}
\begin{SGAresult}{练习}{2.3.6}
（\Term{standard-resolution}{标准消解}。）

a) 设 $C$ 是\Term{small-category}{小范畴}。对每个整数 $n>0$，记 $\operatorname{Fl}^n(C)$ 为 $C$ 中\Term{morphism}{态射}列 $(u_1,\ldots,u_n)$ 的\Term{set}{集合}，要求对每个 $i$，$0<i<n$，$u_i$ 的靶等于 $u_{i+1}$ 的源，从而 $u_i$ 和 $u_{i+1}$ 可复合。
\SourcePage{22}
定义\Term{semisimplicial-set}{半单纯集}
\[
\begin{tikzcd}[column sep=large]
 \operatorname{ES}(C)=\operatorname{ob}C&
 \operatorname{Fl}^1(C)\arrow[l,shift left=1]\arrow[l,shift right=1]&
 \operatorname{Fl}^2(C)\arrow[l,shift left=2]\arrow[l]\arrow[l,shift right=2]&
 \operatorname{Fl}^3(C)\arrow[l,shift left=3]\arrow[l,shift left=1]\arrow[l,shift right=1]\arrow[l,shift right=3]&\cdots,
\end{tikzcd}
\]
其\Term{face-operator}{面算子} $s_i:\operatorname{Fl}^n(C)\to\operatorname{Fl}^{n-1}(C)$ 为
\[
\begin{aligned}
 s_0(u_1,\ldots,u_n)&=(u_2,\ldots,u_n),\\
 s_i(u_1,\ldots,u_n)&=(u_1,\ldots,u_{i+1}u_i,\ldots,u_n),&&0<i<n,\\
 s_n(u_1,\ldots,u_n)&=(u_1,\ldots,u_{n-1}).
\end{aligned}
\]
b) 证明：若 $C$ 有\Term{initial-object}{始对象}或\Term{final-object}{终对象}，则\Term{complex}{复形} $\operatorname{ES}(C)$ \Term{homotopically-trivial}{同伦平凡}。

c) 对 $C$ 的每个\Term{object}{对象} $X$，记 $X\backslash C$ 为源为 $X$ 的\Term{arrow-category}{箭头范畴}。定义一个\Term{semisimple}{半单}纯\Term{presheaf}{预层} $\mathscr{ES}(C)$，使其在每个\Term{object}{对象} $X$ 处的值为 $\operatorname{ES}(X\backslash C)$。记 $\mathbf Z_{\mathscr{ES}(C)}$ 为 $\mathscr{ES}(C)$ 生成的自由阿贝尔\Term{semisimple}{半单}纯\Term{presheaf}{预层}。它带有到值为 $\mathbf Z$ 的常\Term{presheaf}{预层}的\Term{canonical}{典范}\Term{augmentation}{增广} $\mathbf Z_{\mathscr{ES}(C)}\to\mathbf Z$。证明：取相伴的\Term{differential}{微分}\Term{complex}{复形}后，$\mathbf Z_{\mathscr{ES}(C)}$ 是 $\mathbf Z$ 的\Term{resolution}{消解}，且其各项是\Term{projective-module}{投射}阿贝尔\Term{presheaf}{预层}。

d) 对每个阿贝尔\Term{presheaf}{预层} $M$，置 $\operatorname{ST}^\bullet(M)=\operatorname{Hom}(\mathbf Z_{\mathscr{ES}(C)},M)$。写出 $\operatorname{ST}^\bullet(M)$ 的各项。证明对每个整数 $q\geq0$，有
$H^q(\operatorname{ST}^\bullet(M))=H^q(\widehat C,M)$。

e) 证明对 $C$ 上的每个\Term{sheaf}{层} $S$，\Term{functor}{函子} $H^q(S,-)$ 与\Term{restriction-of-scalars}{标量限制}交换。

f) 注意：当 $C$ 是群时，$\mathbf Z_{\mathscr{ES}(C)}$ 是平凡模 $\mathbf Z$ 的\Term{standard-resolution}{标准消解} [3]。

g) 证明对 $C$ 上的每个阿贝尔\Term{presheaf}{预层} $M$，$\mathbf Z_{\mathscr{ES}(C)}\otimes M$ 是 $M$ 的一个（左）\Term{resolution}{消解}。在\Term{presheaf}{预层}上定义\Term{functor}{函子} $\varinjlim_C$，证明它\Term{right-exact}{右正合}。记其左\Term{derived-functor}{导出函子}为 $H_q(C,M)$。证明\Term{complex}{复形} $\mathbf Z_{\mathscr{ES}(C)}\otimes M$ 的各项对 $\varprojlim_C$ \Term{acyclic-object}{无环}。由此，把\Term{complex}{复形} $\varprojlim_C(\mathbf Z_{\mathscr{ES}(C)}\otimes M)$ 记为 $\operatorname{ST}_\bullet(M)$，得出\Term{isomorphism}{同构}
$H_q(\operatorname{ST}_\bullet(M))\simeq H_q(C,M)$。

h) 设 $u:\widehat C\to\operatorname{Ens}$ 是从\Term{topos}{拓扑斯} $\widehat C$ 到点\Term{topos}{拓扑斯}的唯一\Term{morphism}{态射}。记 $u_!:\widehat C_{\mathbf Z}\to\operatorname{Ab}$ 为\Term{functor}{函子} $u^*:\operatorname{Ab}\to\widehat C_{\mathbf Z}$ 的\Term{left-adjoint}{左伴随}。证明对每个阿贝尔\Term{presheaf}{预层} $M$，$u_!M=\varinjlim_CM$。

i) 此后把原来记为 $\operatorname{ST}^\bullet(M)$、$\operatorname{ST}_\bullet(M)$ 的\Term{complex}{复形}分别记为 $\operatorname{ST}^\bullet(C,M)$、$\operatorname{ST}_\bullet(C,M)$。若 $C$ 上的\Term{presheaf}{预层} $M$ 把 $C$ 的每个\Term{morphism}{态射}都变为\Term{isomorphism}{同构}，则称它\Term{locally-constant}{局部常值}。给每个\Term{locally-constant}{局部常值}\Term{presheaf}{预层} $M$ 配上\Term{opposite-category}{反范畴} $C^\circ$ 上的\Term{locally-constant}{局部常值}\Term{presheaf}{预层} $\check M$，使对 $C$ 的每个\Term{morphism}{态射} $u$，有 $\check M(u)=M(u)^{-1}$。求\Term{complex}{复形} $\operatorname{ST}^\bullet(C,M)$ 与 $\operatorname{ST}^\bullet(C^\circ,\check M)$（相应地 $\operatorname{ST}_\bullet(C,M)$ 与 $\operatorname{ST}_\bullet(C^\circ,\check M)$）之间的\Term{canonical-isomorphism}{典范同构}。
\SourcePage{23}
j) 设 $C$ 是有\Term{initial-object}{始对象}的\Term{small-category}{小范畴}（相应地，小\Term{filtered-category}{滤过范畴}）。证明对每个常\Term{presheaf}{预层} $M$，$q>0$ 时有 $H^q(\widehat C,M)=0$（相应地，$q>0$ 时有 $H_q(C,M)=0$）。
\end{SGAresult}

\textbf{2.4。\Term{small-site}{小景}的情形；\Term{cech-cohomology}{Čech 上同调}。}

\textbf{2.4.1。}设 $(C,A)$ 是\Term{ringed}{赋环} $\mathcal U$-\Term{site}{景}，$\widehat C$ 是 $C$ 上的\Term{sheaf}{层}\Term{topos}{拓扑斯}，$\varepsilon:C\to\widehat C$ 是把 $C$ 的\Term{object}{对象}映到其所代表的\Term{presheaf}{预层}的\Term{associated-sheaf}{伴随层}的\Term{canonical-functor}{典范函子}。滥用记号，对 $C$ 的每个\Term{object}{对象} $X$ 和每个 $A$-模\Term{sheaf}{层} $M$，置 $H^q(X,M)=H^q(X,M)$（参见 2.3.1）。回顾一下：当 $C$ 的\Term{topology}{拓扑}比\Term{canonical-topology}{典范拓扑}粗时——这在实际中总是成立——$\varepsilon$ \Term{fully-faithful}{全忠实}，可把 $C$ 与其在 $\varepsilon$ 下的像等同。

\textbf{2.4.2。}记 $\mathscr H^0:\widetilde C_A\to\widehat C_A$ 为从 $A$-模\Term{sheaf}{层}\Term{category}{范畴}到 $A$-模\Term{presheaf}{预层}\Term{category}{范畴}的\Term{inclusion-functor}{包含函子}。因此，对每个 $A$-模\Term{sheaf}{层} $M$，按定义有
\[
 \mathscr H^0(M)(X)=H^0(X,M)=M(X),\tag{2.4.2.1}
\]
其中 $X$ 是 $C$ 的任意\Term{object}{对象}。\Term{functor}{函子} $\mathscr H^0$ \Term{left-exact}{左正合}，其右\Term{derived-functor}{导出函子}记为 $\mathscr H^q$。由于对每个 $X\in C$，“$X$ 上的\Term{section}{截面}”\Term{functor}{函子}在\Term{presheaf}{预层}\Term{category}{范畴}中\Term{exact}{正合}，有
\[
 \mathscr H^q(M)(X)=H^q(X,M),\tag{2.4.2.2}
\]
其中 $X$ 为 $C$ 的任意\Term{object}{对象}，$M$ 为任意 $A$-模\Term{sheaf}{层}。故\Term{presheaf}{预层} $\mathscr H^q(M)$ 正是 $X\mapsto H^q(X,M)$。

\textbf{2.4.3。}假设 $(C,A)$ 是小的\Term{ringed}{赋环}\Term{site}{景}，使 $\widehat C$ 是可应用 2.3 的结果的\Term{topos}{拓扑斯}。设 $X\in C$，$R\hookrightarrow X$ 是\Term{covering-sieve}{覆盖筛}。对每个 $A$-模\Term{presheaf}{预层} $G$，群 $H^q(R,G)$（在\Term{topos}{拓扑斯} $\widehat C$ 中计算）称为\Term{presheaf}{预层} $G$ 相对于\Term{covering-sieve}{覆盖筛} $R$ 的\Term{cech-cohomology}{Čech 上同调群}。若 $R\hookrightarrow X$ 是由\Term{quarrable}{可作基变换}\Term{morphism}{态射}的\Term{covering-family}{覆盖族} $\mathcal U=(U_i\to X)_{i\in I}$ 生成的\Term{sieve}{筛}，则这些群可用\Term{complex}{复形} $C^\bullet(\mathcal U,G)$（2.3.3）计算，称其为 $G$ 相对于\Term{covering-family}{覆盖族} $\mathcal U$（或\Term{covering-recouvrement}{覆盖} $\mathcal U$）的\Term{cech-complex}{Čech 复形}。群 $H^q(\mathcal U,G)=H^q(R,G)$（2.3.4）称为 $G$ 相对于\Term{covering-family}{覆盖族} $\mathcal U$ 的\Term{cech-cohomology}{Čech 上同调群}。

\textbf{2.4.4。}设 $M$ 是 $C$ 上的 $A$-模\Term{sheaf}{层}。群 $H^q(\mathcal U,\mathscr H^0(M))$ 常被滥记为 $H^q(\mathcal U,M)$，并称为\Term{sheaf}{层} $M$ 相对于\Term{covering-family}{覆盖族} $\mathcal U$ 的\Term{cech-cohomology}{Čech 上同调群}。

\textbf{2.4.5。}记 $\check{\mathscr H}^0:\widehat C_A\to\widehat C_A$ 为第 II 讲所述\Term{functor}{函子} $L$ 到 $A$-模\Term{presheaf}{预层}的自然推广。因此，对\Term{presheaf}{预层} $G$ 及 $C$ 的\Term{object}{对象} $X$，按定义有
\SourcePage{24}
\[
 \check{\mathscr H}^0(G)(X)=\varinjlim_{R\hookrightarrow X}G(R),\tag{2.4.5.1}
\]
其中\Term{inductive-limit}{归纳极限}沿\Term{covering}{覆盖} $X$ 的\Term{sieve}{筛}取。由 (2.4.5.1)，\Term{functor}{函子} $\check{\mathscr H}^0$ \Term{left-exact}{左正合}，其右\Term{derived-functor}{导出函子}记为 $\check{\mathscr H}^q$。由于“$X$ 上的\Term{section}{截面}”和“\Term{filtered-inductive-limit}{滤过归纳极限}”\Term{functor}{函子}\Term{exact}{正合}，由 (2.4.5.1) 得
\[
 \check{\mathscr H}^q(G)(X)=\varinjlim_{R\hookrightarrow X}H^q(R,G),\tag{2.4.5.2}
\]
其中\Term{limit}{极限}沿\Term{covering}{覆盖} $X$ 的\Term{sieve}{筛}取。\Term{presheaf}{预层} $\check{\mathscr H}^q$ 称为\Term{cech-cohomology}{Čech 上同调}\Term{presheaf}{预层}。对 $C$ 的每个\Term{object}{对象} $X$，置
\[
 \check H^q(X,G)=\check{\mathscr H}^q(G)(X).\tag{2.4.5.3}
\]
群 $\check H^q(X,G)$ 称为\Term{cech-cohomology}{Čech 上同调群}。当 $C$ 的\Term{topology}{拓扑}由\Term{pretopology}{预拓扑}定义时——这在实际中最为常见——由 2.3.4 得
\[
 \check H^q(X,G)\simeq\varinjlim_{\mathcal U}H^q(\mathcal U,G),\tag{2.4.5.4}
\]
其中\Term{inductive-limit}{归纳极限}沿\Term{quarrable}{可作基变换}的\Term{covering-family}{覆盖族}取，它们由对应\Term{sieve}{筛}之间的自然\Term{order-relation}{序关系}预序化（2.3.5）。

\textbf{2.4.6。}设 $M$ 是 $A$-模\Term{sheaf}{层}。滥用记号，置
\[
 \check H^q(X,M)=\check H^q(X,\mathscr H^0(M)),\qquad
 \check{\mathscr H}^q(M)=\check{\mathscr H}^q(\mathscr H^0(M)).\tag{2.4.6.1}
\]
群 $\check H^q(X,M)$ 称为\Term{sheaf}{层} $M$ 的\Term{cech-cohomology}{Čech 上同调群}。须指出：尽管\Term{functor}{函子} $\check H^q$ 在\Term{presheaf}{预层}\Term{category}{范畴}上是\Term{derived-functor}{导出函子}，它们一般并不在\Term{sheaf}{层}\Term{category}{范畴}上组成 $\delta$-\Term{functor}{函子}。

\textbf{2.5。\Term{universe}{宇宙}变换；$\mathcal U$-景情形的 \Term{cech-cohomology}{Čech 上同调}。}

\textbf{2.5.1。}设 $(C,A)$ 是\Term{ringed}{赋环} $\mathcal U$-\Term{site}{景}，$\mathcal V$ 是包含 $\mathcal U$ 的\Term{universe}{宇宙}。于是 $(C,A)$ 是 $\mathcal V$-\Term{site}{景}，有 $\mathcal U$-\Term{topos}{拓扑斯} $\widetilde C_{\mathcal U}$、$\mathcal V$-\Term{topos}{拓扑斯} $\widetilde C_{\mathcal V}$ 以及\Term{canonical}{典范}\Term{inclusion-functor}{包含函子} $\varepsilon:\widetilde C_{\mathcal U}\hookrightarrow\widetilde C_{\mathcal V}$。在\Term{module-category}{模范畴}上，$\varepsilon$ \Term{exact}{正合}且\Term{fully-faithful}{全忠实}，并把\Term{injective-object}{内射模}变为\Term{injective-object}{内射模}（1.9）。因此，对每对 $A$-模 $\mathcal U$-\Term{sheaf}{层}，有\Term{canonical-isomorphism}{典范同构}
\[
 \operatorname{Ext}_A^q(\widetilde C_{\mathcal U};M,N)
 \simeq\operatorname{Ext}_{\varepsilon A}^q(\widetilde C_{\mathcal V};\varepsilon M,\varepsilon N),
 \qquad q\geq0.\tag{2.5.1.1}
\]
特别地，对 $C$ 上的每个 $\mathcal U$-\Term{set}{集合}\Term{sheaf}{层} $R$，有\Term{canonical-isomorphism}{典范同构}（2.1.1）
\[
 H^q(R,M)\simeq H^q(\varepsilon R,\varepsilon M),\qquad q\geq0,\tag{2.5.1.2}
\]
\SourcePage{25}
更特别地，对 $C$ 的每个\Term{object}{对象} $X$，有\Term{canonical-isomorphism}{典范同构}（2.4.1）
\[
 H^q(X,M)\simeq H^q(X,\varepsilon M).\tag{2.5.1.3}
\]
粗略地说，\Term{sheaf-cohomology}{层上同调}不依赖\Term{universe}{宇宙}的选择；为证明或构造之需，总可以扩大\Term{universe}{宇宙}来计算\Term{sheaf}{层}的\Term{cohomology}{上同调}。

\textbf{2.5.2。}设 $(C,A)$ 是\Term{ringed}{赋环} $\mathcal U$-\Term{site}{景}，$\mathcal V$ 是包含 $\mathcal U$ 的\Term{universe}{宇宙}。记 $\widetilde C_A$ 为 $\mathcal U$-\Term{sheaf}{层}\Term{category}{范畴}，$\widehat\varepsilon:\widehat C_{A,\mathcal U}\to\widehat C_{A,\mathcal V}$ 为从 $A$-模 $\mathcal U$-\Term{presheaf}{预层}到 $\mathcal V$-\Term{presheaf}{预层}的\Term{inclusion-functor}{包含函子}。$\widehat\varepsilon$ \Term{exact}{正合}，故\Term{functor}{函子}
$\widehat\varepsilon\mathscr H^0:\widetilde C_A\to\widehat C_{A,\mathcal V}$
的\Term{derived-functor}{导出函子}是 $\widehat\varepsilon\mathscr H^q$，$q\geq0$。仍滥用记号，把 $\widehat\varepsilon\mathscr H^q$ 记为
$\mathscr H^q:\widetilde C_A\to\widehat C_{A,\mathcal V}$，$q\geq0$。
当 $C$ 是 $\mathcal V$-小时，扩大\Term{universe}{宇宙}有如下好处：\Term{category}{范畴} $\widehat C_{\mathcal U}$ 一般不是 $\mathcal U$-\Term{topos}{拓扑斯}，$A$-模 $\mathcal U$-\Term{presheaf}{预层}也不一定是\Term{injective-object}{内射} $\mathcal U$-\Term{presheaf}{预层}的\Term{submodule}{子模}；但 $\mathcal V$-\Term{presheaf}{预层}\Term{category}{范畴}是\Term{topos}{拓扑斯}（$\mathcal V$-\Term{topos}{拓扑斯}），所以每个 $A$-模 $\mathcal V$-\Term{presheaf}{预层}都是\Term{injective-object}{内射} $\mathcal V$-\Term{presheaf}{预层}的\Term{subobject}{子对象}（0.1.1）。因此，对每个 $\mathcal V$-\Term{set}{集合}\Term{presheaf}{预层} $R$（特别地，$R$ 为 $\mathcal U$-\Term{presheaf}{预层}时），及每个 $A$-模 $\mathcal U$-\Term{sheaf}{层} $M$，群 $H^q(R,\mathscr H^q(M))$ 由 2.1.1 定义；由 (2.5.1.2)，这些群不依赖所考虑的\Term{universe}{宇宙} $\mathcal V$。同样，对每对非负整数 $p,q$，\Term{presheaf}{预层} $\check{\mathscr H}^p(\mathscr H^q(M))$ 由 2.4.5 定义；由 (2.5.1.2) 和 (2.4.5.4)，这些\Term{presheaf}{预层}不依赖用于定义它们的\Term{universe}{宇宙} $\mathcal V$。

\SGAsection{3}{与覆盖有关的 Cartan--Leray 谱序列}
\begin{SGAresult}{命题}{3.1}
设 $(C,A)$ 是\Term{ringed}{赋环} $\mathcal U$-\Term{site}{景}，$\mathcal V$ 是包含 $\mathcal U$ 的\Term{universe}{宇宙}。\Term{functor}{函子} $\mathscr H^0:\widetilde C_A\to\widehat C_{A,\mathcal V}$（2.5.2）把\Term{injective-object}{内射} $A$-模变为\Term{injective-object}{内射}\Term{presheaf}{预层}。对每个整数 $q>0$ 及每个 $A$-模 $M$，\Term{presheaf}{预层} $\mathscr H^q(M)$ 的\Term{associated-sheaf}{伴随层}为零。
\end{SGAresult}
记 $a_{\mathcal V}$ 为 $\mathcal V$-\Term{presheaf}{预层}的“\Term{associated-sheaf}{伴随层}”\Term{functor}{函子}，$\varepsilon:\widetilde C_A\hookrightarrow\widetilde C_{A,\mathcal V}$ 为 $\mathcal U$-\Term{sheaf}{层}到 $\mathcal V$-\Term{sheaf}{层}的\Term{inclusion-functor}{包含函子}。有 $a_{\mathcal V}\mathscr H^0=\varepsilon$（第 II 讲 3.6）；由于 $a_{\mathcal V}$ 与 $\varepsilon$ \Term{exact}{正合}，$q>0$ 时有 $a_{\mathcal V}\mathscr H^q=0$。对每个 $\mathcal U$-\Term{sheaf}{层} $M$ 和每个 $\mathcal V$-\Term{presheaf}{预层} $N$，有\Term{functor}{函子}\Term{isomorphism}{同构}
\[
 \operatorname{Hom}_{\widehat C_{A,\mathcal V}}(N,\mathscr H^0(M))
 \simeq\operatorname{Hom}_{\widetilde C_{A,\mathcal V}}(a_{\mathcal V}N,\varepsilon M).
\]
若 $M$ \Term{injective-object}{内射}，则 $\varepsilon M$ \Term{injective-object}{内射}（1.9）；由于 $a_{\mathcal V}$ \Term{exact}{正合}，\Term{functor}{函子} $N\mapsto\operatorname{Hom}_{\widehat C_{A,\mathcal V}}(N,\mathscr H^0(M))$ \Term{exact}{正合}。因此 $\mathscr H^0(M)$ \Term{injective-object}{内射}。
\begin{SGAresult}{定理}{3.2}
设 $(C,A)$ 是\Term{ringed}{赋环} $\mathcal U$-\Term{site}{景}，$R$ 是 $C$ 上的 $\mathcal U$-\Term{set}{集合}\Term{presheaf}{预层}，$M$ 是 $A$-模\Term{sheaf}{层}。存在关于 $R$、$M$ \Term{functorial}{函子性}的\Term{spectral-sequence}{谱序列}
\SourcePage{26}
\[
 H^{p+q}(aR,M)\Longleftarrow E_2^{p,q}=H^p(R,\mathscr H^q(M)).\tag{3.2.1}
\]
这里 $aR$ 表示 $R$ 的\Term{associated-sheaf}{伴随层}。
（当 $C$ 不是 $\mathcal U$-小时，$H^p(R,\mathscr H^q(M))$ 应解释为\Term{presheaf}{预层} $\mathscr H^q(M)$ 在\Term{topos}{拓扑斯} $\widehat C_{\mathcal V}$ 中的\Term{cohomology}{上同调}，其中 $\mathcal V$ 包含 $\mathcal U$，且使 $C$ 为 $\mathcal V$-小（2.5.2）。）
\end{SGAresult}
由“\Term{associated-sheaf}{伴随层}”\Term{functor}{函子}的定义（参见第 II 讲），有\Term{functor}{函子}\Term{isomorphism}{同构} $H^0(aR,M)\simeq H^0(R,\mathscr H^0(M))$。$\mathscr H^0$ 把\Term{injective-object}{内射对象}变为\Term{injective-object}{内射对象}。\Term{composite-functor}{复合函子}的\Term{spectral-sequence}{谱序列}（0.3）就是所求的\Term{spectral-sequence}{谱序列}。
\begin{SGAresult}{推论}{3.3}
设 $X\in C$，$\mathcal U=(U_i\to X)_{i\in I}$ 是\Term{covering-family}{覆盖族}，且每个\Term{morphism}{态射} $U_i\to X$ 都\Term{quarrable}{可作基变换}。则有（称为 Cartan--Leray 的）\Term{spectral-sequence}{谱序列}
\[
 H^{p+q}(X,M)\Longleftarrow E_2^{p,q}=H^p(\mathcal U,\mathscr H^q(M)).\tag{3.3.1}
\]
\end{SGAresult}
设 $R\hookrightarrow X$ 是 $\mathcal U$ 生成的\Term{sieve}{筛}。由于此\Term{sieve}{筛}为\Term{covering-sieve}{覆盖筛}，$R$ 的\Term{associated-sheaf}{伴随层}就是 $X$ 的\Term{associated-sheaf}{伴随层}（第 II 讲 5.2）。由 2.4.1，$H^{p+q}(aR,M)=H^{p+q}(X,M)$。于是由 2.3.4 得到本推论。
\begin{SGAresult}{推论}{3.4}
存在关于\Term{sheaf}{层} $M$ 和\Term{object}{对象} $X\in C$ \Term{functorial}{函子性}的\Term{spectral-sequence}{谱序列}
\[
 H^{p+q}(X,M)\Longleftarrow E_2^{p,q}=\check H^p(X,\mathscr H^q(M)).\tag{3.4.1}
\]
当 $X$ 在 $C$ 中变化时，它给出\Term{presheaf}{预层}的\Term{spectral-sequence}{谱序列}
\[
 \mathscr H^{p+q}(M)\Longleftarrow E_2^{p,q}=\check{\mathscr H}^p(\mathscr H^q(M)).\tag{3.4.2}
\]
这些\Term{spectral-sequence}{谱序列}给出\Term{functorial}{函子性}\Term{morphism}{态射}（\Term{edge-morphism}{边缘态射}）
\[
 \Phi^q(M):\check{\mathscr H}^q(M)\longrightarrow\mathscr H^q(M),\tag{3.4.3}
\]
\[
 \Phi_X^q(M):\check H^q(X,M)\longrightarrow H^q(X,M).\tag{3.4.4}
\]
$q=0$ 或 $1$ 时，$\Phi^q$、$\Phi_X^q$ 是\Term{isomorphism}{同构}；$\Phi^2$、$\Phi_X^2$ 是\Term{monomorphism}{单态射}。更一般地，若 $0<i<n$ 时\Term{presheaf}{预层} $\mathscr H^i(M)$ 为零，则 $0\leq q\leq n$ 时 $\Phi^q(M)$、$\Phi_X^q(M)$ 为\Term{isomorphism}{同构}，$q=n+1$ 时为\Term{monomorphism}{单态射}。
\end{SGAresult}
对\Term{covering}{覆盖} $X$ 的\Term{sieve}{筛} $R\hookrightarrow X$，在 (3.2.1) 中取\Term{inductive-limit}{归纳极限}即得第一个\Term{spectral-sequence}{谱序列}（2.4.5.2）。第二个立即由第一个得到（2.4.5.3）。$q>0$ 时，\Term{presheaf}{预层} $\mathscr H^q(M)$ 的\Term{associated-sheaf}{伴随层}为零（3.1）。故 $q>0$ 时
$\check{\mathscr H}^0\check{\mathscr H}^0\mathscr H^q(M)=0$（第 II 讲 3.4）。
\SourcePage{27}
于是 $\check{\mathscr H}^0\mathscr H^q(M)$ 是其\Term{associated-sheaf}{伴随层}为零的\Term{separated}{分离}\Term{presheaf}{预层}。因此 $q>0$ 时 $\check{\mathscr H}^0\mathscr H^q(M)=0$（第 II 讲 3.2）。关于\Term{morphism}{态射} $\Phi^q$、$\Phi_X^q$ 的断言立即由此得到。
\begin{SGAresult}{推论}{3.5}
设 $(E,A)$ 是\Term{ringed-topos}{赋环拓扑斯}，$M$ 是 $A$-模（为明确起见，取左模）。记 $\mathcal M$ 为 $M$ 的底层\Term{abelian-group}{阿贝尔群}\Term{object}{对象}。\Term{functor}{函子} $M\mapsto\mathcal M$ \Term{exact}{正合}；因此，对 $E$ 的每个\Term{object}{对象} $X$，\Term{canonical-isomorphism}{典范同构}
\[
 H^0(X,M)\xrightarrow{\sim}H^0(X,\mathcal M)
\]
延拓为\Term{morphism}{态射}
\[
 H^q(X,M)\longrightarrow H^q(X,\mathcal M),\qquad q\geq0.
\]
这些\Term{morphism}{态射}都是\Term{isomorphism}{同构}。
\end{SGAresult}
对 $E$ 的每个\Term{object}{对象} $Y$，有（2.4.5.4）
\[
 H^p(Y,M)=\varinjlim_{\mathcal U}H^p(\mathcal U,M),\qquad
 H^p(Y,\mathcal M)=\varinjlim_{\mathcal U}H^p(\mathcal U,\mathcal M),
\]
其中\Term{inductive-limit}{归纳极限}沿\Term{epimorphic-family}{满态射族} $\mathcal U=(Y_i\to Y)_{i\in I}$ 取。因此\Term{canonical-homomorphism}{典范同态} $\check H^p(Y,M)\to\check H^p(Y,\mathcal M)$ 是\Term{isomorphism}{同构}（2.3.4）；故由 (2.4.5.3)，\Term{canonical-homomorphism}{典范同态} $\check{\mathscr H}^p(M)\to\check{\mathscr H}^p(\mathcal M)$ 是\Term{isomorphism}{同构}。若 $M$ 是\Term{injective-object}{内射} $A$-模，则 $p>0$ 时 $\check{\mathscr H}^p(\mathcal M)=0$；从而 $\mathscr H^1(\mathcal M)=0$，对 $p$ 归纳，得到 $p>0$ 时 $\mathscr H^p(\mathcal M)=0$（3.4）。所以 $p>0$ 时 $H^p(X,\mathcal M)=0$（2.4.2.2 及 2.5）。因此\Term{functor}{函子} $M\mapsto\mathcal M$ 把\Term{injective-object}{内射对象}变为对 $H^0(X,-)$ \Term{acyclic-object}{无环的对象}，从而得到所述\Term{isomorphism}{同构}。
\begin{SGAresult}{练习}{3.6}
设 $G$ 是\Term{topological-group}{拓扑群}，$B_G$ 为其\Term{classifying-topos}{分类拓扑斯}（第 IV 讲 2.5）。记 $E_G$ 为 $B_G$ 的如下\Term{object}{对象}：$G$ 的底层\Term{topological-space}{拓扑空间}，带有 $G$ 的元素的\Term{left-translation}{左平移}\Term{action}{作用}。从 $E_G$ 到 $B_G$ 的\Term{final-object}{终对象} $e_G$ 的\Term{canonical}{典范}\Term{morphism}{态射}是\Term{epimorphism}{满态射}，故给出\Term{covering-recouvrement}{覆盖} $\mathcal U=(E_G\to e_G)$；对 $B_G$ 的每个阿贝尔\Term{sheaf}{层} $F$，得到\Term{spectral-sequence}{谱序列}
\[
 E_2^{p,q}=H^p(\mathcal U,\mathscr H^q(F))\Longrightarrow H^{p+q}(B_G,F),\tag{3.6.1}
\]
我们拟研究此\Term{spectral-sequence}{谱序列}。

1) 证明对每个整数 $n$，\Term{topos}{拓扑斯}
\[
 B_G/(E_G\times E_G\times\cdots\times E_G)\quad\text{（$n$ 个因子）}
\]
\Term{canonical}{典范}\Term{equivalence}{等价}于\Term{topos}{拓扑斯}（第 IV 讲 2.5）
\[
 \operatorname{TOP}(G\times G\times\cdots\times G)\quad\text{（$n-1$ 个因子）}.
\]
对每个整数 $n$，记 $F_n$ 为 $F$ 所诱导的\Term{topological-space}{拓扑空间} $G\times\cdots\times G$（$n$ 个因子）上的\Term{sheaf}{层}。
\SourcePage{28}
2) 由此推得，(3.6.1) 的 $E_2^{p,q}$ 项\Term{canonical-isomorphism}{典范同构}于如下类型的\Term{complex}{复形}的第 $p$ 个\Term{cohomology-group}{上同调群}：
\[
 H^q(e,F_0)\longrightarrow H^q(G,F_1)\longrightarrow
 H^q(G\times G,F_2)\longrightarrow\cdots,
\]
并写出此\Term{complex}{复形}（$e$ 表示单点\Term{topological-space}{拓扑空间}）。

3) 由此推得：如果对于 $G\times G\times\cdots\times G$ 类型的\Term{topological-space}{拓扑空间}，\Term{locally-constant}{局部常值}\Term{sheaf}{层}的\Term{cohomology}{上同调}与相应的\Term{singular-cohomology}{奇异上同调}一致，那么\Term{classifying-topos}{分类拓扑斯} $B_G$ 以\Term{locally-constant}{局部常值}\Term{sheaf}{层}为系数的\Term{cohomology}{上同调}，\Term{isomorphism}{同构}于\Term{topology}{拓扑}学中使用的分类\Term{space}{空间}的相应\Term{singular-cohomology}{奇异上同调}。（例如 $G$ \Term{local-scheme}{局部}可缩时，上述前提成立；关于分类\Term{space}{空间}，可参阅 [15]。）
\end{SGAresult}

\SGAsection{4}{无环层}
\begin{SGAresult}{定义}{4.1}
设 $(E,A)$ 是\Term{ringed-topos}{赋环拓扑斯}，$F$ 是 $A$-模，$S$ 是 $E$ 的一个\Term{topologically-generating-family}{拓扑生成族}。若对 $S$ 的每个\Term{object}{对象} $X$ 和每个整数 $q>0$，有 $H^q(X,F)=0$，则称 $F$ 是 $S$-\Term{acyclic-sheaf}{无环的}。当 $S=\operatorname{ob}E$ 时，$S$-\Term{acyclic-sheaf}{无环层}称为\Term{flasque-sheaf}{松层}。
\end{SGAresult}
\textbf{4.2。}设 $(C,A)$ 是\Term{ringed}{赋环} $\mathcal U$-\Term{site}{景}，$F$ 是 $C$ 上的 $A$-模\Term{sheaf}{层}。若 $F$ 是 $S$-\Term{acyclic-sheaf}{无环的}，其中 $S$ 是 $C$ 的\Term{object}{对象}的\Term{associated-sheaf}{伴随层}组成的族，则称 $F$ 为 $C$-\Term{acyclic-sheaf}{无环的}。
\begin{SGAresult}{命题}{4.3}
设 $(C,A)$ 是\Term{ringed}{赋环} $\mathcal U$-\Term{site}{景}，$F$ 是 $A$-模\Term{sheaf}{层}。记 $\mathscr H^0:\widetilde C_A\to\widehat C_A$ 为从 $A$-\Term{module-category}{模范畴}到 $A$-模\Term{presheaf}{预层}\Term{category}{范畴}的\Term{inclusion-functor}{包含函子}（这里的\Term{presheaf}{预层}是 $\mathcal V$-\Term{presheaf}{预层}，$\mathcal V$ 是使 $C$ 为 $\mathcal V$-小的\Term{universe}{宇宙}）。下列条件\Term{equivalence}{等价}：

\noindent i) $F$ 是 $C$-\Term{acyclic-sheaf}{无环的}。

\noindent ii) 对 $C$ 的每个\Term{object}{对象} $X$ 和每个\Term{covering-sieve}{覆盖筛} $R\hookrightarrow X$，有
\[
 H^q(R,\mathscr H^0(F))=0\qquad\text{对每个 $q>0$}.
\]
\noindent iii) 对 $C$ 的每个\Term{object}{对象} $X$，有
\[
 \check H^q(X,F)=0\qquad\text{对每个 $q>0$}.
\]
\end{SGAresult}
$\mathrm{i})\Rightarrow\mathrm{ii})$：由于 $F$ 是 $C$-\Term{acyclic-sheaf}{无环的}，$q>0$ 时\Term{presheaf}{预层} $\mathscr H^q(F)$ 为零。\Term{spectral-sequence}{谱序列} (3.2.1) 对每个 $q$ 给出\Term{isomorphism}{同构} $H^q(R,\mathscr H^0(F))\simeq H^q(X,F)$。故 $q>0$ 时 $H^q(R,\mathscr H^0(F))=0$。

$\mathrm{ii})\Rightarrow\mathrm{iii})$：取\Term{inductive-limit}{归纳极限}即可。

$\mathrm{iii})\Rightarrow\mathrm{i})$：利用 (3.4.2)，对 $q$ 归纳证明 $q>0$ 时\Term{presheaf}{预层} $\mathscr H^q(F)$ 为零。

\SourcePage{29}
\textbf{4.4。}由判据 4.3 ii) 和 3.5，$S$-\Term{acyclic-sheaf}{无环性}只依赖底层阿贝尔\Term{sheaf}{层}。特别地，$A$-模\Term{sheaf}{层}为\Term{flasque-sheaf}{松层}，当且仅当其底层阿贝尔\Term{sheaf}{层}为\Term{flasque-sheaf}{松层}。
\begin{SGAresult}{推论}{4.5}
设 $(E,A)$ 是\Term{ringed-topos}{赋环拓扑斯}，$F$ 是 $A$-模。下列性质\Term{equivalence}{等价}：

\noindent i) $F$ 为\Term{flasque-sheaf}{松层}；

\noindent ii) 对每个\Term{epimorphic-family}{满态射族} $\mathcal X=(X_i\to X)_{i\in I}$，有
\[
 H^q(\mathcal X,F)=0\qquad\text{对每个 $q>0$}.
\]
\end{SGAresult}
在 4.2 中取 $C$ 为\Term{topos}{拓扑斯} $E$ 本身。推论由 4.2 的\Term{equivalence}{等价}性 $\mathrm{i})\Leftrightarrow\mathrm{ii})$ 得到。

\textbf{4.6。}\Term{injective-object}{内射}\Term{sheaf}{层}都是\Term{flasque-sheaf}{松层}。对每个\Term{topologically-generating-family}{拓扑生成族} $S$，\Term{flasque-sheaf}{松层}都是 $S$-\Term{acyclic-sheaf}{无环的}。\Term{flasque-sheaf}{松层}不一定\Term{injective-object}{内射}（取 $E$ 为\Term{set}{集合}\Term{topos}{拓扑斯}）。$S$-\Term{acyclic-sheaf}{无环层}不一定为\Term{flasque-sheaf}{松层}（练习 4.1.3）。
\begin{SGAresult}{命题}{4.7}
设 $(E,A)$ 是\Term{topos}{拓扑斯}，$F$ 是一个\Term{flasque-sheaf}{松} $A$-模，$X$ 是 $E$ 的\Term{object}{对象}。对 $X$ 的每个\Term{subobject}{子对象} $Y$，\Term{canonical-homomorphism}{典范同态} $H^0(X,F)\to H^0(Y,F)$ 为\Term{surjectivity}{满射}。
\end{SGAresult}
设 $Y$ 是 $X$ 的一个\Term{subobject}{子对象}，使 $H^0(X,F)\to H^0(Y,F)$ 不是\Term{surjectivity}{满射}。设 $Z$ 是把 $X$ 的两个副本 $X_1,X_2$ 沿 $Y$ \Term{gluing}{黏合}所得的\Term{object}{对象}。$Z$ 被两个\Term{subobject}{子对象} $X_1,X_2$ \Term{covering}{覆盖}，且 $X_1\times_ZX_2=Y$。记所得\Term{covering-recouvrement}{覆盖}为 $\mathcal X=(X_1,X_2)$。立即可见 $H^1(\mathcal X,F)\ne0$，矛盾。

\textbf{4.8。}在一般\Term{topos}{拓扑斯}中，4.7 所述性质并不刻画\Term{flasque-sheaf}{松层}（练习 4.15）。但在由其\Term{open-object}{开对象}生成的\Term{topos}{拓扑斯}中，特别是在\Term{topological-space}{拓扑空间}所对应的\Term{topos}{拓扑斯}中，它确实刻画\Term{flasque-sheaf}{松层}（练习 4.16）。这里采用的\Term{flasque-sheaf}{松层}术语在\Term{topological-space}{拓扑空间}的情形下与 [7] 的术语一致（练习 4.16）。
\begin{SGAresult}{命题}{4.9}
设 $u:(E,A)\to(E',A')$ 是\Term{ringed-topos}{赋环拓扑斯}\Term{morphism}{态射}。

1) \Term{functor}{函子} $u_*$ 把\Term{flasque-sheaf}{松} $A$-模变为\Term{flasque-sheaf}{松} $A'$-模。

2) 设 $S,S'$ 分别是 $E,E'$ 的\Term{topologically-generating-family}{拓扑生成族}，且 $u^*(S')\subset S$。\Term{functor}{函子} $u_*$ 把 $S$-\Term{acyclic-sheaf}{无环} $A$-模变为 $S'$-\Term{acyclic-sheaf}{无环} $A'$-模。

3) 若 $u$ 是\Term{flat-morphism}{平坦态射}（1.8），则 $u_*$ 把\Term{injective-object}{内射} $A$-模变为\Term{injective-object}{内射} $A'$-模。
\end{SGAresult}
\SourcePage{30}
设 $X$ 是 $E'$ 的\Term{object}{对象}，$\mathcal X=(X_i\to X)_{i\in I}$ 是\Term{epimorphic-family}{满态射族}，$F$ 是\Term{flasque-sheaf}{松} $A$-模，$C^\bullet(\mathcal X,u_*F)$ 是\Term{covering-recouvrement}{覆盖} $\mathcal X$ 的\Term{cech-complex}{Čech 复形}。利用 $u_*$ 与 $u^*$ 的\Term{adjunction}{伴随}性以及 $u^*$ 与\Term{fiber-product}{纤维积}交换，有\Term{canonical-isomorphism}{典范同构}
\[
 C^\bullet(\mathcal X,u_*F)\simeq C^\bullet(u^*(\mathcal X),F).
\]
此外，$u^*$ 与\Term{inductive-limit}{归纳极限}交换，所以 $u^*(\mathcal X)=(u^*X_i\to u^*X)_{i\in I}$ 是\Term{epimorphic-family}{满态射族}。由于 $F$ 是\Term{flasque-sheaf}{松层}，$q>0$ 时 $H^q(u^*(\mathcal X),F)=0$，故
\[
\begin{split}
 H^q(\mathcal X,u_*F)&=H^q(C^\bullet(\mathcal X,u_*F))\\
 &=H^q(C^\bullet(u^*(\mathcal X),F))=H^q(u^*(\mathcal X),F)=0
 \qquad(q>0).
\end{split}
\]
因此 $u_*F$ 是\Term{flasque-sheaf}{松层}。

证明 2)。若 $S'$ 对\Term{fiber-product}{纤维积}\Term{stable}{稳定}，可用与上面类似的证明得到 2)。一般情形下，将使用\Term{morphism}{态射} $u$ 的\Term{spectral-sequence}{谱序列}（5.3.2）；在 5.3.2 之前不会使用断言 2)。设 $F$ 是 $S$-\Term{acyclic-sheaf}{无环层}。\Term{sheaf}{层} $R^qu_*F$ 是\Term{presheaf}{预层} $X\mapsto H^q(u^*X,F)$ 的\Term{associated-sheaf}{伴随层}。由于 $S'$ 是\Term{topologically-generating-family}{拓扑生成族}，且 $F$ 为 $S$-\Term{acyclic-sheaf}{无环的}，对每个 $q>0$ 有 $R^qu_*F=0$。于是\Term{spectral-sequence}{谱序列} (5.3.2) 对每个 $X\in E'$ 给出\Term{isomorphism}{同构} $H^q(X,u_*F)\simeq H^q(u^*X,F)$；从而对每个 $X\in S'$，$H^q(X,u_*F)=0$。

证明 3)。若 $u$ \Term{flat-morphism}{平坦}，模上的\Term{functor}{函子} $u^*$ \Term{exact}{正合}（1.8）。所以 $u^*$ 的\Term{right-adjoint}{右伴随} $u_*$ 把\Term{injective-object}{内射对象}变为\Term{injective-object}{内射对象}（0.2）。
\begin{SGAresult}{命题}{4.10}
设 $F$ 是\Term{ringed-topos}{赋环拓扑斯} $(E,A)$ 的 $A$-模，$G$ 是\Term{injective-object}{内射} $A$-模。

1) \Term{functor}{函子} $F\mapsto\mathscr{H}om_A(F,G)$ \Term{exact}{正合}。

2) \Term{abelian-group}{阿贝尔群}\Term{object}{对象} $\mathscr{H}om_A(F,G)$ 为\Term{flasque-sheaf}{松层}。
\end{SGAresult}
\begin{proof}
1) 设 $0\to F'\to F\to F''\to0$ 是\Term{exact-sequence}{正合列}。须证明
\[
 0\longrightarrow\mathscr{H}om_A(F'',G)\longrightarrow\mathscr{H}om_A(F,G)
 \longrightarrow\mathscr{H}om_A(F',G)\longrightarrow0
\]
\Term{exact}{正合}；为此只须对 $E$ 的每个\Term{object}{对象} $H$，证明下列序列\Term{exact}{正合}：
\[
\begin{split}
 0\longrightarrow&\operatorname{Hom}_E(H,\mathscr{H}om_A(F'',G))
 \longrightarrow\operatorname{Hom}_E(H,\mathscr{H}om_A(F,G))\\
 &\longrightarrow\operatorname{Hom}_E(H,\mathscr{H}om_A(F',G))\longrightarrow0.
\end{split}
\]
由第 IV 讲 6.12，最后这个序列\Term{isomorphism}{同构}于
\[
\begin{split}
 0\longrightarrow&\operatorname{Hom}_A(A_H\otimes_AF'',G)
 \longrightarrow\operatorname{Hom}_A(A_H\otimes_AF,G)\\
 &\longrightarrow\operatorname{Hom}_A(A_H\otimes_AF',G)\longrightarrow0,
\end{split}
\]
而 $A$-模 $A_H$ \Term{flat-module}{平坦}（1.3.1）。

2) 证明 $\mathscr{H}om_A(F,G)$ 是\Term{flasque-sheaf}{松层}。设 $\mathcal X=(X_i\to X)_{i\in I}$
\SourcePage{31}
是 $E$ 的\Term{epimorphic-family}{满态射族}，并设
\[
\begin{tikzcd}[column sep=large]
 C_\bullet(\mathcal X):\cdots\arrow[r,shift left=2]\arrow[r]\arrow[r,shift right=2]&
 \displaystyle\coprod_{i,j}\mathbf Z_{X_i\times_X X_j}
 \arrow[r,shift left=1]\arrow[r,shift right=1]&\displaystyle\coprod_i\mathbf Z_{X_i}
\end{tikzcd}
\]
为 1.5 定义的\Term{complex}{复形}。它是\Term{object}{对象} $\mathbf Z_X$ 的\Term{flat-resolution}{平坦消解}，而 $\mathbf Z_X$ 本身\Term{flat-module}{平坦}（1.4 和 1.3）。计算群：
\[
\begin{split}
 H^q(\mathcal X,\mathscr{H}om_A(F,G))
 &\xrightarrow{\sim}H^q(\operatorname{Hom}_{\mathbf Z}^{\bullet}(C_\bullet(\mathcal X),\mathscr{H}om_A(F,G)))\\
 &\xrightarrow{\sim}H^q(\operatorname{Hom}_A(C_\bullet(\mathcal X)\otimes_{\mathbf Z}F,G)).
\end{split}
\]
\Term{complex}{复形} $C_\bullet(\mathcal X)\otimes_{\mathbf Z}F$ 在非零次数上\Term{acyclic-complex}{无环}，因为 $C_\bullet(\mathcal X)$ 是\Term{flat-module}{平坦模}的\Term{flat-resolution}{平坦消解}。所以 $\operatorname{Hom}(C_\bullet(\mathcal X)\otimes_{\mathbf Z}F,G)$ 在非零次数上\Term{acyclic-complex}{无环}。
\end{proof}
\begin{SGAresult}{命题}{4.11}
设 $(E,A)$ 是\Term{ringed-topos}{赋环拓扑斯}，$F$ 是松（相应地\Term{injective-object}{内射}）$A$-模\Term{sheaf}{层}。

1) 对 $E$ 的每个\Term{object}{对象} $X$，$A|X$-模 $j_X^*F$ 是松的（相应地\Term{injective-object}{内射的}）。

2) 对 $E$ 的每个闭集 $Z$，$F$ 中\Term{support}{支集}包含于 $Z$ 的\Term{section}{截面}\Term{sheaf}{层}（第 IV 讲 14）为\Term{flasque-sheaf}{松的}（相应地\Term{injective-object}{内射的}）。
\end{SGAresult}
$F$ 为\Term{flasque-sheaf}{松层}时，断言 1) 由 4.5 得到。$j_X^*$ 有\Term{exact}{正合}的\Term{left-adjoint}{左伴随} $j_{X!}$（第 IV 讲 11），故 $F$ \Term{injective-object}{内射}时 $j_X^*F$ \Term{injective-object}{内射}（0.2）。证明 2)。记 $i:Z\to E$ 为\Term{inclusion-morphism}{包含态射}。$F$ 中\Term{support}{支集}包含于 $Z$ 的\Term{section}{截面}\Term{sheaf}{层}此时是 $i_*i^!F$（第 IV 讲 14）。\Term{functor}{函子} $i_*i^!$ 是\Term{exact-functor}{正合函子} $i_*i^*$ 的\Term{right-adjoint}{右伴随}（第 IV 讲 14），故把\Term{injective-object}{内射}\Term{sheaf}{层}变为\Term{injective-object}{内射}\Term{sheaf}{层}。设 $U$ 是 $Z$ 的补开子\Term{topos}{拓扑斯}，$j:U\to E$ 是\Term{inclusion-morphism}{包含态射}，$F$ 为\Term{flasque-sheaf}{松层}。有\Term{exact-sequence}{正合列}（第 IV 讲 14）
\[
 0\longrightarrow i_*i^!F\longrightarrow F\longrightarrow j_*j^*F.\tag{4.11.1}
\]
对 $E$ 的每个\Term{object}{对象} $X$，有 $j_*j^*F(X)=F(X\times U)$，由 (4.11.1) 最后一个\Term{arrow}{箭头}诱导的\Term{morphism}{态射} $F(X)\to j_*j^*F(X)$ 来自\Term{canonical}{典范}\Term{injectivity}{单射} $X\times U\hookrightarrow X$。由于 $F$ 是\Term{flasque-sheaf}{松层}，此\Term{morphism}{态射}为\Term{surjectivity}{满射}（4.7），故 (4.11.1) 的最后一个\Term{arrow}{箭头}为\Term{presheaf}{预层}的\Term{epimorphism}{满态射}。对 $E$ 的每个\Term{object}{对象} $X$，由 (4.11.1) \Term{derivation}{导出}的\Term{cohomology}{上同调}\Term{exact-sequence}{正合列}给出 $q>0$ 时 $H^q(X,i_*i^!F)=0$，故 $i_*i^!F$ 为\Term{flasque-sheaf}{松层}。

\textbf{4.12。}\Term{sheaf}{层}为\Term{flasque-sheaf}{松的}（相应地\Term{injective-object}{内射的}）这一性质可\Term{localization}{局部化}（4.11.1），但它通常不是\Term{local-property}{局部性质}（练习 4.15）。然而，在由其\Term{open-object}{开对象}生成的\Term{topos}{拓扑斯}中，特别是在\Term{topological-space}{拓扑空间}所对应的\Term{topos}{拓扑斯}中，它是\Term{local-property}{局部性质}（练习 4.16）。
\SourcePage{32}
\begin{SGAresult}{练习}{4.13}
设 $X$ 是\Term{locally-compact-space}{局部紧空间}，$F$ 是 $c$-\Term{soft-sheaf}{软层} [7]。利用软性的\Term{local-property}{局部性}（同上），证明 $F$ 在 $X$ 的每个仿紧\Term{open-subset}{开集}上的\Term{restriction}{限制}都是\Term{soft-sheaf}{软层}。证明仿紧\Term{open-subset}{开集}构成 $X$ 的\Term{topology}{拓扑}的一组基。由此推得：若记 $S$ 为 $X$ 的仿紧\Term{open-subset}{开集}族，则 $F$ 是 $S$-\Term{acyclic-sheaf}{无环的}。证明 $X$ 上的\Term{continuous}{连续}函数\Term{sheaf}{层}是 $S$-\Term{acyclic-sheaf}{无环的}，但当 $X$ 非\Term{discrete}{离散}时不是\Term{flasque-sheaf}{松层}。
\end{SGAresult}
\begin{SGAresult}{问题}{4.14}
研究\Term{totally-acyclic-topos}{全无环拓扑斯}，即对每个 $q>0$、每个\Term{object}{对象} $X$ 和每个阿贝尔\Term{sheaf}{层} $F$ 都有 $H^q(X,F)=0$ 的\Term{topos}{拓扑斯}。
\end{SGAresult}
\begin{SGAresult}{练习}{4.15}
设 $G$ 是\Term{discrete}{离散}群，$B_G$ 是其\Term{classifying-topos}{分类拓扑斯}（第 IV 讲 2.4）。证明对每个阿贝尔\Term{sheaf}{层} $F$ 和每个\Term{monomorphism}{单态射} $X\hookrightarrow Y$，\Term{homomorphism}{同态} $F(Y)\to F(X)$ 为\Term{surjectivity}{满射}。设 $E(G)$ 是把群 $G$ 看成其自身的\Term{homogeneous-space}{齐性空间}所得的 $B_G$ 的\Term{object}{对象}。\Term{topos}{拓扑斯} $B_G/E(G)$ \Term{equivalence}{等价}于点\Term{topos}{拓扑斯}（第 IV 讲 8）。\Term{morphism}{态射} $E(G)\to e$（$e$ 为 $B_G$ 的\Term{final-object}{终对象}）是\Term{epimorphism}{满态射}。对 $B_G$ 的每个阿贝尔\Term{sheaf}{层} $F$，$F|E(G)$ 是\Term{flasque-sheaf}{松层}。由此推得：\Term{flasque-sheaf}{松性}或\Term{injective-object}{内射性}不是\Term{local-property}{局部性质}。
\end{SGAresult}
\begin{SGAresult}{练习}{4.16}
若\Term{topos}{拓扑斯} $E$ 的\Term{open-object}{开对象}（即 $E$ 的\Term{final-object}{终对象} $e$ 的\Term{subobject}{子对象}）构成一个\Term{generating-family}{生成族}（第 I 讲 7），则称 $E$ 由其\Term{open-object}{开对象}生成。这种\Term{topos}{拓扑斯}有以下性质：

(P) 每个\Term{epimorphic-family}{满态射族} $(X_i\to X)_{i\in I}$ 都以某个\Term{epimorphic-family}{满态射族} $(U_j\to X)_{j\in J}$ 为上界，其中 $U_j\to X$ 都是\Term{monomorphism}{单态射}。

a) 是否存在满足 (P) 而不由其\Term{open-object}{开对象}生成的\Term{topos}{拓扑斯}？\Term{topological-space}{拓扑空间}所对应的\Term{topos}{拓扑斯}由其\Term{open-object}{开对象}生成。若 $E$ 由其\Term{open-object}{开对象}生成（相应地满足 (P)），则对 $E$ 的每个\Term{object}{对象} $X$，$E/X$ 由其\Term{open-object}{开对象}生成（相应地满足 (P)）。由其\Term{open-object}{开对象}生成的\Term{topos}{拓扑斯}的每个子\Term{topos}{拓扑斯}也由其\Term{open-object}{开对象}生成。性质 (P) 不是\Term{local-property}{局部性质}。

b) 设 $E$ 是\Term{topos}{拓扑斯}，$\operatorname{Ouv}(E)$ 是 $E$ 的\Term{open-object}{开对象}\Term{category}{范畴}，赋予\Term{induced-topology}{诱导拓扑}。\Term{inclusion-functor}{包含函子} $\operatorname{Ouv}(E)\to E$ 是\Term{site}{景}\Term{morphism}{态射}，故给出\Term{topos}{拓扑斯}\Term{morphism}{态射} $\Pi:E\to\operatorname{Ouv}(E)^\sim$。\Term{functor}{函子} $\Pi^*$ \Term{fully-faithful}{全忠实}。\Term{morphism}{态射} $\Pi$ 相对于由其\Term{open-object}{开对象}生成的\Term{topos}{拓扑斯}的 $2$-\Term{category}{范畴}具有一个\Term{universal-property}{泛性质}，请读者明确写出。

c) 设 $E$ 是由其\Term{open-object}{开对象}生成的\Term{topos}{拓扑斯}，$\operatorname{Point}(E)$ 是 $E$ 的点按\Term{isomorphism}{同构}分类所得的\Term{set}{集合}。证明 $\operatorname{Point}(E)$ 是小的。在 $\operatorname{Point}(E)$ 上赋予\Term{topology}{拓扑}，并定义\Term{morphism}{态射}
\[
 \operatorname{Top}(\operatorname{Point}(E))\longrightarrow\operatorname{Ouv}(E)^\sim,
\]
使 $\operatorname{Top}(\operatorname{Point}(E))$ 成为 $\operatorname{Ouv}(E)^\sim$ 的一个子\Term{topos}{拓扑斯}。对每个\Term{topological-space}{拓扑空间} $X$，证明从 $\operatorname{Top}(X)$ 到 $E$ 的每个\Term{morphism}{态射}都给出从 $\operatorname{Top}(X)$ 到 $\operatorname{Top}(\operatorname{Point}(E))$ 的\Term{morphism}{态射}。证明由其\Term{open-object}{开对象}生成的\Term{topos}{拓扑斯}\Term{equivalence}{等价}于某个\Term{topos}{拓扑斯} $\operatorname{Top}(X)$（$X$ 为\Term{topological-space}{拓扑空间}），当且仅当它有\Term{enough-points}{足够多的点}。
\SourcePage{33}
证明存在由其\Term{open-object}{开对象}生成、但不\Term{equivalence}{等价}于任何\Term{topological-space}{拓扑空间} $X$ 对应的 $\operatorname{Top}(X)$ 的\Term{topos}{拓扑斯}。

d) 设 $E$ 是满足 (P) 的\Term{topos}{拓扑斯}。$E$ 上的阿贝尔\Term{sheaf}{层} $F$ 为\Term{flasque-sheaf}{松层}的充要条件是：对每个\Term{monomorphism}{单态射} $X\hookrightarrow Y$，\Term{homomorphism}{同态} $F(Y)\to F(X)$ 为\Term{surjectivity}{满射}。

e) 设 $E$ 是满足 (P) 的\Term{topos}{拓扑斯}。证明\Term{sheaf}{层} $F$ 的\Term{flasque-sheaf}{松性}是\Term{local-property}{局部性质}。

f) 设 $E$ 是由其\Term{open-object}{开对象}生成的\Term{topos}{拓扑斯}，$e$ 是 $E$ 的\Term{final-object}{终对象}。阿贝尔\Term{sheaf}{层} $F$ 为\Term{flasque-sheaf}{松层}的充要条件是：对 $e$ 的每个\Term{open-object}{开对象} $U$，\Term{canonical}{典范}\Term{morphism}{态射} $F(e)\to F(U)$ 为\Term{surjectivity}{满射}。
\end{SGAresult}
\begin{SGAresult}{练习}{4.17}
（\Term{flasque-sheaf}{松层}与\Term{universe}{宇宙}变换。）设 $C$ 是 $\mathcal U$-\Term{site}{景}（例如 $\mathcal U$-\Term{topos}{拓扑斯}），$\mathcal V$ 是包含 $\mathcal U$ 的\Term{universe}{宇宙}。记相应的\Term{sheaf}{层}\Term{category}{范畴}及\Term{canonical}{典范}包含为 $\varepsilon:\widetilde C_{\mathcal U}\to\widetilde C_{\mathcal V}$。设 $F$ 是 $C$ 上的\Term{flasque-sheaf}{松}阿贝尔 $\mathcal U$-\Term{sheaf}{层}。拟证明 $\varepsilon F$ 为\Term{flasque-sheaf}{松层}。首先注意对 $\widetilde C_{\mathcal U}$ 的每个\Term{object}{对象} $X$，$q>0$ 时 $H^q(\varepsilon X,\varepsilon F)=0$。$\widetilde C_{\mathcal V}$ 的每个\Term{object}{对象} $Y$ 有一个\Term{epimorphic-family}{满态射族} $\varepsilon X_i\to Y$，$i\in I$，其中 $\varepsilon X_i\to Y$ 都是\Term{monomorphism}{单态射}。由此得出
\[
 H^q(Y,\varepsilon F)=\varprojlim_{\varepsilon X\hookrightarrow Y}^{q}F(X).
\]
要证明这些 $\varprojlim^{q}$ 在 $q\ne0$ 时为零，可归结到\Term{topos}{拓扑斯} $\operatorname{Ouv}(\widetilde C_{\mathcal V}/Y)^\sim$，并利用如下事实：在这个\Term{topos}{拓扑斯}中，\Term{local-scheme}{局部}为\Term{flasque-sheaf}{松的}\Term{sheaf}{层}就是\Term{flasque-sheaf}{松层}。
\end{SGAresult}
% 本块至此结束；第5节标题及以下归下一块。
